\clearpage
\fancyhead[L]{\small NADIRPUNKT.}
\SourceLeaf{477}
\begin{center}
\rule{0.55\linewidth}{0.4pt}\\[2pt]
Astronomische Nachrichten. Band XXV. Nr. 579. Seite 43. 44. 1846 December 10.\\[2pt]
\rule{0.55\linewidth}{0.4pt}
\end{center}
\begin{flushright}Göttingen 1846 October 31.\end{flushright}

Wenn man den Nadirpunkt durch das vom Quecksilberhorizonte reflectirte Bild der Fäden bestimmen will, so wünscht man namentlich das Gesichtsfeld einigermassen gleichförmig beleuchtet zu haben. Das war aber früher bei dem Gebrauche einer Lampe nie der Fall. Es zeigte sich vielmehr nur ein etwas undeutliches Bild der Lichtflamme in dem sonst äusserst schwach erhellten Felde.

Ich habe früher mir dadurch geholfen, dass ich das Ocular etwas seitwärts schraubte, so dass nicht der mittelste Faden in der Mitte, sondern seitwärts von der Mitte erschien, etwa so, dass der nächste Faden in der Mitte des Gesichtsfeldes stand. Dadurch erlangte ich den Vortheil, dass jenes Flammenbild die Beobachtungen nicht mehr bedeutend störte. Es ging so besser, als wenn das Ocular concentrisch stand, aber immer noch nicht so gut wie man wünschen konnte. Es war dies also nur ein Palliativ. Radical wird auf folgende Art geholfen.

\begin{center}
\begingroup
\setlength{\unitlength}{1bp}
\begin{picture}(228,195)
\put(0,0){\special{pdf:bcontent}\special{pdf:code q 0 g 56.700000 178.800000 m 56.400000 178.500000 l 54.300000 178.500000 l 54.000000 178.200000 l 53.400000 178.200000 l 53.100000 177.900000 l 52.800000 178.200000 l 52.500000 177.900000 l 51.900000 177.900000 l 51.600000 178.200000 l 51.300000 178.200000 l 51.000000 177.900000 l 50.400000 177.900000 l 50.100000 177.600000 l 49.200000 177.600000 l 48.900000 177.300000 l 47.700000 177.300000 l 47.400000 177.000000 l 46.800000 177.000000 l 46.500000 176.700000 l 45.900000 176.700000 l 45.600000 176.400000 l 45.300000 176.400000 l 45.000000 176.100000 l 44.400000 176.100000 l 44.400000 175.800000 l 44.100000 175.500000 l 43.800000 175.500000 l 43.500000 175.200000 l 43.200000 175.200000 l 42.900000 175.500000 l 42.300000 174.900000 l 41.700000 174.900000 l 40.500000 173.700000 l 40.200000 173.700000 l 39.600000 173.100000 l 39.600000 172.800000 l 39.300000 172.500000 l 39.300000 172.200000 l 39.000000 171.900000 l 39.000000 171.600000 l 38.700000 171.300000 l 35.100000 171.300000 l 34.500000 171.900000 l 34.200000 171.600000 l 31.800000 171.600000 l 31.500000 171.300000 l 30.600000 171.300000 l 30.300000 171.600000 l 26.100000 171.600000 l 25.800000 171.900000 l 22.200000 171.900000 l 22.200000 171.600000 l 21.900000 171.300000 l 22.200000 171.000000 l 22.200000 170.700000 l 22.500000 170.400000 l 22.500000 169.800000 l 22.200000 169.500000 l 22.200000 166.800000 l 22.500000 166.500000 l 22.500000 163.500000 l 22.200000 163.200000 l 22.500000 162.900000 l 22.500000 161.100000 l 22.200000 160.800000 l 22.200000 154.800000 l 22.500000 154.500000 l 22.500000 154.200000 l 22.200000 153.900000 l 22.500000 153.600000 l 22.200000 153.300000 l 22.200000 153.000000 l 21.900000 152.700000 l 22.200000 152.400000 l 22.200000 150.300000 l 21.900000 150.000000 l 22.200000 149.700000 l 22.200000 149.400000 l 21.900000 149.100000 l 21.900000 148.500000 l 22.200000 148.200000 l 21.900000 147.900000 l 21.900000 147.300000 l 22.200000 147.000000 l 22.200000 146.700000 l 21.900000 146.400000 l 22.200000 146.100000 l 22.200000 144.600000 l 21.900000 144.300000 l 21.900000 142.500000 l 22.200000 142.200000 l 22.200000 140.400000 l 22.500000 140.100000 l 22.200000 139.800000 l 22.200000 136.200000 l 21.900000 135.900000 l 21.900000 132.600000 l 22.200000 132.300000 l 21.900000 132.000000 l 21.900000 131.400000 l 22.200000 131.100000 l 21.900000 130.800000 l 21.900000 130.500000 l 22.200000 130.200000 l 22.200000 128.400000 l 22.500000 128.100000 l 21.900000 127.500000 l 22.200000 127.200000 l 21.900000 126.900000 l 21.900000 125.700000 l 21.600000 125.400000 l 21.300000 125.400000 l 21.900000 125.400000 l 22.200000 125.100000 l 22.200000 123.900000 l 21.900000 123.600000 l 22.200000 123.300000 l 22.200000 120.900000 l 21.900000 120.600000 l 22.200000 120.300000 l 22.200000 119.400000 l 21.900000 119.100000 l 21.900000 118.500000 l 22.200000 118.200000 l 21.900000 117.900000 l 22.200000 117.600000 l 22.200000 116.100000 l 22.500000 115.800000 l 22.200000 115.500000 l 22.200000 112.200000 l 22.500000 111.900000 l 22.200000 111.600000 l 22.200000 111.000000 l 21.900000 110.700000 l 21.900000 108.900000 l 21.600000 108.600000 l 21.900000 108.300000 l 21.900000 106.200000 l 22.200000 105.900000 l 22.200000 104.700000 l 21.900000 104.400000 l 22.200000 104.100000 l 21.900000 103.800000 l 22.200000 103.500000 l 22.200000 102.600000 l 21.900000 102.300000 l 21.900000 99.600000 l 21.600000 99.300000 l 21.900000 99.000000 l 21.900000 97.800000 l 21.600000 97.500000 l 21.900000 97.200000 l 21.900000 96.900000 l 21.600000 96.600000 l 21.900000 96.300000 l 21.900000 95.100000 l 21.600000 94.800000 l 21.600000 94.500000 l 21.900000 94.200000 l 21.600000 93.900000 l 21.600000 93.300000 l 21.900000 93.000000 l 21.900000 92.700000 l 21.600000 92.400000 l 21.900000 92.100000 l 21.900000 91.500000 l 21.600000 91.200000 l 21.900000 90.900000 l 21.900000 90.300000 l 21.600000 90.000000 l 21.900000 89.700000 l 21.900000 88.500000 l 21.600000 88.200000 l 21.600000 87.600000 l 21.900000 87.300000 l 21.900000 87.000000 l 21.600000 86.700000 l 21.600000 86.400000 l 21.900000 86.100000 l 21.300000 85.500000 l 21.300000 85.200000 l 21.600000 84.900000 l 21.600000 84.300000 l 21.300000 84.000000 l 21.300000 83.400000 l 21.600000 83.100000 l 21.300000 82.800000 l 21.300000 82.200000 l 21.600000 81.900000 l 21.600000 81.600000 l 21.900000 81.300000 l 21.900000 81.000000 l 21.600000 80.700000 l 21.600000 79.800000 l 21.300000 79.500000 l 21.600000 79.200000 l 21.600000 78.900000 l 21.900000 78.600000 l 21.600000 78.300000 l 21.600000 77.400000 l 21.300000 77.100000 l 21.300000 75.300000 l 21.000000 75.000000 l 21.300000 74.700000 l 21.300000 69.900000 l 21.600000 69.600000 l 21.600000 69.000000 l 21.300000 68.700000 l 21.300000 64.800000 l 21.600000 64.500000 l 21.300000 64.200000 l 21.300000 62.700000 l 21.600000 62.400000 l 21.600000 62.100000 l 21.300000 61.800000 l 21.300000 61.500000 l 21.600000 61.200000 l 21.300000 60.900000 l 21.600000 60.600000 l 21.300000 60.300000 l 21.300000 60.000000 l 21.000000 59.700000 l 21.000000 58.800000 l 21.300000 58.500000 l 21.000000 58.200000 l 21.600000 57.600000 l 21.600000 55.800000 l 21.300000 55.500000 l 21.600000 55.200000 l 21.600000 52.800000 l 21.000000 52.200000 l 21.000000 51.900000 l 21.300000 51.600000 l 21.300000 51.000000 l 21.000000 50.700000 l 21.300000 50.400000 l 21.300000 48.300000 l 21.000000 48.000000 l 21.300000 47.700000 l 21.300000 47.100000 l 21.000000 46.800000 l 21.000000 46.500000 l 21.300000 46.500000 l 21.600000 46.200000 l 22.800000 46.200000 l 23.100000 46.500000 l 23.400000 46.500000 l 23.700000 46.200000 l 24.000000 46.500000 l 24.300000 46.200000 l 24.300000 45.900000 l 24.600000 46.200000 l 24.900000 45.900000 l 25.200000 46.200000 l 25.500000 46.200000 l 25.800000 45.900000 l 26.100000 46.200000 l 27.600000 46.200000 l 27.900000 45.900000 l 28.500000 45.900000 l 28.800000 46.200000 l 29.100000 46.200000 l 29.400000 46.500000 l 29.700000 46.200000 l 31.500000 46.200000 l 31.800000 46.500000 l 32.400000 46.500000 l 32.700000 46.200000 l 36.900000 46.200000 l 37.200000 46.500000 l 37.800000 45.900000 l 37.800000 45.300000 l 38.100000 45.300000 l 38.400000 45.000000 l 38.400000 44.700000 l 39.000000 44.100000 l 39.000000 43.800000 l 39.300000 43.500000 l 39.300000 43.200000 l 39.600000 42.900000 l 39.900000 42.900000 l 40.800000 42.000000 l 40.800000 41.700000 l 41.400000 41.700000 l 41.700000 41.400000 l 42.300000 41.400000 l 42.600000 41.100000 l 42.600000 40.800000 l 42.900000 41.100000 l 43.200000 40.800000 l 43.500000 40.800000 l 43.800000 40.500000 l 44.400000 40.500000 l 45.000000 39.900000 l 45.300000 40.200000 l 45.600000 39.900000 l 46.200000 39.900000 l 46.500000 39.600000 l 46.800000 39.600000 l 47.100000 39.300000 l 48.000000 39.300000 l 48.300000 39.000000 l 48.900000 39.000000 l 49.200000 38.700000 l 49.500000 38.700000 l 49.800000 38.400000 l 51.000000 38.400000 l 51.300000 38.100000 l 53.100000 38.100000 l 53.400000 38.400000 l 53.700000 38.400000 l 54.000000 38.100000 l 55.800000 38.100000 l 56.100000 37.800000 l 56.400000 37.800000 l 56.700000 38.100000 l 57.300000 38.100000 l 57.600000 37.800000 l 57.600000 37.500000 l 59.400000 37.500000 l 59.700000 37.200000 l 60.000000 37.500000 l 65.400000 37.500000 l 65.700000 37.800000 l 66.000000 37.500000 l 68.100000 37.500000 l 68.400000 37.800000 l 69.900000 37.800000 l 70.200000 38.100000 l 71.700000 38.100000 l 72.000000 38.400000 l 72.300000 38.100000 l 72.600000 38.400000 l 72.900000 38.400000 l 73.200000 38.100000 l 73.500000 38.400000 l 73.800000 38.400000 l 74.100000 38.700000 l 74.700000 38.100000 l 75.000000 38.400000 l 75.300000 38.400000 l 75.600000 38.700000 l 76.500000 38.700000 l 76.800000 39.000000 l 77.400000 39.000000 l 78.000000 39.600000 l 78.900000 39.600000 l 79.200000 39.900000 l 79.800000 39.900000 l 79.800000 40.200000 l 80.100000 40.500000 l 81.000000 40.500000 l 81.600000 41.100000 l 81.900000 41.100000 l 82.200000 41.400000 l 82.800000 41.400000 l 83.100000 41.700000 l 83.400000 41.700000 l 84.300000 42.600000 l 84.600000 42.600000 l 84.900000 42.900000 l 85.200000 42.900000 l 85.200000 43.500000 l 85.500000 43.800000 l 85.500000 44.400000 l 85.800000 44.700000 l 85.800000 45.000000 l 86.100000 45.300000 l 86.400000 45.300000 l 86.700000 45.600000 l 87.000000 45.300000 l 89.700000 45.300000 l 90.000000 45.600000 l 90.900000 45.600000 l 91.200000 45.300000 l 91.200000 45.000000 l 91.500000 45.300000 l 91.800000 45.000000 l 92.100000 45.300000 l 92.400000 45.000000 l 93.600000 45.000000 l 93.900000 45.300000 l 94.200000 45.000000 l 96.300000 45.000000 l 96.600000 45.300000 l 96.900000 45.000000 l 97.800000 45.000000 l 98.100000 45.300000 l 98.400000 45.000000 l 99.300000 45.000000 l 99.600000 44.700000 l 101.700000 44.700000 l 102.000000 45.000000 l 104.100000 45.000000 l 104.400000 45.300000 l 105.000000 45.300000 l 105.300000 45.600000 l 105.300000 49.500000 l 105.600000 49.800000 l 105.300000 50.100000 l 105.900000 50.700000 l 105.600000 51.000000 l 105.900000 51.300000 l 105.900000 51.900000 l 105.600000 52.200000 l 105.900000 52.500000 l 105.600000 52.800000 l 105.900000 53.100000 l 105.600000 53.400000 l 105.900000 53.700000 l 105.600000 54.000000 l 105.600000 54.900000 l 105.900000 55.200000 l 105.900000 55.800000 l 105.600000 56.100000 l 105.600000 58.500000 l 105.300000 58.800000 l 105.600000 59.100000 l 105.900000 59.100000 l 105.900000 62.100000 l 105.600000 62.400000 l 105.600000 62.700000 l 105.900000 63.000000 l 105.900000 63.900000 l 106.200000 64.200000 l 106.200000 64.500000 l 105.900000 64.800000 l 106.200000 65.100000 l 105.900000 65.400000 l 105.900000 66.300000 l 105.600000 66.600000 l 105.900000 66.900000 l 105.900000 67.200000 l 106.200000 67.500000 l 106.200000 67.800000 l 105.900000 68.100000 l 105.900000 69.300000 l 106.200000 69.600000 l 105.600000 70.200000 l 106.200000 70.800000 l 105.900000 71.100000 l 105.900000 71.700000 l 105.600000 72.000000 l 105.900000 72.300000 l 105.900000 72.600000 l 106.200000 72.900000 l 106.200000 73.800000 l 105.900000 74.100000 l 105.900000 74.700000 l 106.200000 75.000000 l 106.200000 75.300000 l 105.900000 75.600000 l 105.900000 75.900000 l 106.500000 76.500000 l 106.500000 77.100000 l 106.800000 77.400000 l 106.500000 77.400000 l 106.200000 77.700000 l 106.500000 78.000000 l 106.500000 78.900000 l 106.200000 79.200000 l 106.500000 79.500000 l 106.500000 79.800000 l 106.200000 80.100000 l 106.500000 80.400000 l 106.500000 80.700000 l 106.200000 81.000000 l 106.500000 81.300000 l 106.200000 81.600000 l 106.200000 83.400000 l 105.900000 83.700000 l 106.500000 84.300000 l 106.500000 85.200000 l 106.800000 85.500000 l 106.200000 86.100000 l 106.500000 86.400000 l 106.200000 86.700000 l 106.500000 87.000000 l 106.500000 89.400000 l 106.200000 89.700000 l 106.800000 90.300000 l 106.500000 90.600000 l 106.500000 92.100000 l 106.200000 92.100000 l 105.900000 92.400000 l 106.200000 92.700000 l 106.200000 93.000000 l 106.500000 93.300000 l 106.500000 93.600000 l 106.200000 93.900000 l 106.500000 94.200000 l 106.500000 94.500000 l 106.200000 94.800000 l 106.500000 95.100000 l 106.200000 95.400000 l 106.200000 97.800000 l 106.500000 98.100000 l 107.100000 98.100000 l 107.400000 97.800000 l 107.700000 97.800000 l 108.000000 98.100000 l 108.600000 97.500000 l 109.500000 97.500000 l 109.800000 97.800000 l 109.800000 98.400000 l 109.500000 98.700000 l 109.200000 98.700000 l 108.900000 99.000000 l 106.500000 99.000000 l 105.900000 99.600000 l 106.200000 99.900000 l 106.200000 100.200000 l 105.900000 100.500000 l 106.500000 101.100000 l 106.200000 101.400000 l 106.500000 101.700000 l 106.500000 104.100000 l 106.200000 104.400000 l 106.200000 109.200000 l 106.500000 109.500000 l 106.500000 109.800000 l 106.200000 110.100000 l 106.500000 110.400000 l 106.500000 113.100000 l 106.800000 113.400000 l 106.800000 114.000000 l 106.500000 114.300000 l 106.500000 116.100000 l 106.200000 116.400000 l 106.200000 116.700000 l 106.500000 117.000000 l 106.500000 117.900000 l 106.200000 118.200000 l 106.500000 118.500000 l 106.200000 118.800000 l 106.500000 119.100000 l 106.200000 119.400000 l 106.500000 119.700000 l 106.500000 120.900000 l 106.200000 121.200000 l 106.500000 121.500000 l 106.500000 122.100000 l 106.800000 122.400000 l 106.500000 122.700000 l 106.500000 124.800000 l 106.200000 125.100000 l 106.200000 125.400000 l 106.500000 125.700000 l 106.500000 126.000000 l 106.800000 126.300000 l 107.100000 126.300000 l 106.800000 126.600000 l 106.800000 127.500000 l 107.100000 127.800000 l 106.800000 128.100000 l 106.800000 128.400000 l 106.500000 128.700000 l 106.800000 129.000000 l 106.500000 129.300000 l 106.500000 129.600000 l 106.800000 129.900000 l 106.500000 130.200000 l 106.500000 131.400000 l 106.800000 131.700000 l 106.500000 132.000000 l 106.500000 132.900000 l 106.200000 133.200000 l 106.800000 133.800000 l 106.500000 134.100000 l 106.500000 134.400000 l 106.800000 134.700000 l 106.800000 135.000000 l 107.100000 135.300000 l 107.100000 135.600000 l 106.500000 136.200000 l 107.100000 136.800000 l 107.100000 137.400000 l 106.800000 137.700000 l 106.800000 138.600000 l 107.100000 138.900000 l 107.100000 139.800000 l 106.800000 140.100000 l 106.800000 140.700000 l 106.500000 141.000000 l 106.500000 141.300000 l 106.800000 141.600000 l 106.800000 142.200000 l 106.500000 142.500000 l 106.500000 142.800000 l 106.800000 143.100000 l 107.100000 143.100000 l 107.400000 143.400000 l 107.400000 144.000000 l 107.100000 144.300000 l 107.100000 144.600000 l 107.400000 144.900000 l 106.800000 145.500000 l 107.100000 145.800000 l 106.800000 146.100000 l 106.800000 146.700000 l 107.100000 147.000000 l 106.800000 147.300000 l 106.800000 147.600000 l 107.100000 147.900000 l 108.000000 147.900000 l 108.000000 149.100000 l 107.400000 149.100000 l 106.800000 149.700000 l 107.100000 150.000000 l 107.100000 151.200000 l 107.400000 151.500000 l 107.700000 151.500000 l 107.700000 152.400000 l 107.100000 153.000000 l 107.100000 153.600000 l 107.400000 153.900000 l 107.100000 154.200000 l 107.100000 154.500000 l 107.400000 154.800000 l 107.400000 157.800000 l 107.100000 158.100000 l 107.100000 158.400000 l 107.400000 158.700000 l 107.400000 159.300000 l 107.100000 159.600000 l 107.400000 159.900000 l 107.700000 159.900000 l 107.400000 160.200000 l 107.400000 161.100000 l 107.700000 161.400000 l 107.700000 161.700000 l 107.400000 162.000000 l 107.700000 162.300000 l 107.700000 162.900000 l 107.400000 163.200000 l 107.700000 163.500000 l 107.400000 163.800000 l 107.400000 164.700000 l 107.100000 164.700000 l 106.800000 165.000000 l 107.100000 165.300000 l 107.400000 165.300000 l 107.700000 165.600000 l 107.100000 166.200000 l 107.400000 166.500000 l 107.400000 167.100000 l 107.100000 167.400000 l 107.400000 167.700000 l 107.400000 168.600000 l 107.700000 168.900000 l 107.700000 170.400000 l 105.000000 170.400000 l 104.700000 170.100000 l 103.500000 170.100000 l 103.200000 169.800000 l 102.900000 170.100000 l 102.000000 170.100000 l 101.700000 170.400000 l 101.100000 170.400000 l 100.800000 170.700000 l 100.500000 170.700000 l 100.200000 170.400000 l 97.800000 170.400000 l 97.500000 170.100000 l 97.200000 170.400000 l 93.600000 170.400000 l 93.300000 170.700000 l 92.700000 170.700000 l 92.400000 170.400000 l 88.800000 170.400000 l 87.600000 171.600000 l 87.600000 171.900000 l 87.000000 172.500000 l 86.700000 172.500000 l 86.100000 173.100000 l 86.100000 173.400000 l 85.500000 174.000000 l 85.200000 174.000000 l 84.900000 174.300000 l 84.900000 174.600000 l 84.600000 174.600000 l 84.300000 174.900000 l 84.000000 174.900000 l 83.700000 175.200000 l 83.400000 175.200000 l 83.100000 175.500000 l 82.800000 175.500000 l 82.500000 175.800000 l 81.900000 175.800000 l 81.600000 176.100000 l 81.000000 176.100000 l 80.700000 176.400000 l 80.400000 176.100000 l 80.100000 176.100000 l 79.800000 176.400000 l 78.600000 176.400000 l 78.300000 176.700000 l 77.700000 176.700000 l 77.400000 177.000000 l 76.800000 177.000000 l 76.200000 177.600000 l 75.300000 177.600000 l 75.000000 177.900000 l 73.500000 177.900000 l 73.200000 178.200000 l 72.900000 177.900000 l 72.600000 178.200000 l 71.700000 178.200000 l 71.400000 178.500000 l 71.100000 178.200000 l 70.800000 178.500000 l 68.700000 178.500000 l 68.400000 178.200000 l 68.100000 178.500000 l 66.900000 178.500000 l 66.600000 178.800000 l 66.300000 178.500000 l 63.900000 178.500000 l 63.600000 178.800000 l 62.400000 178.800000 l 62.100000 178.500000 l 61.800000 178.800000 l 61.200000 178.800000 l 60.900000 178.500000 l 60.300000 178.500000 l 60.000000 178.800000 l h 37.500000 46.200000 m 37.800000 46.500000 l 38.100000 46.200000 l 37.800000 45.900000 l h 38.700000 48.000000 m 39.000000 48.300000 l 39.300000 48.000000 l 39.000000 47.700000 l h 54.300000 53.700000 m 54.600000 53.400000 l 54.900000 53.700000 l 55.200000 53.700000 l 55.500000 53.400000 l 55.800000 53.700000 l 56.100000 53.400000 l 56.400000 53.700000 l 56.700000 53.400000 l 57.000000 53.400000 l 57.300000 53.700000 l 57.600000 53.700000 l 57.900000 53.400000 l 58.200000 53.400000 l 58.500000 53.100000 l 58.800000 53.400000 l 65.100000 53.400000 l 65.400000 53.700000 l 65.700000 53.400000 l 66.000000 53.700000 l 66.300000 53.400000 l 66.600000 53.400000 l 66.900000 53.100000 l 69.300000 53.100000 l 69.600000 52.800000 l 69.900000 53.100000 l 70.200000 52.800000 l 72.600000 52.800000 l 72.900000 52.500000 l 73.800000 52.500000 l 74.100000 52.200000 l 74.400000 52.500000 l 75.000000 51.900000 l 75.900000 51.900000 l 76.200000 51.600000 l 77.400000 51.600000 l 77.700000 51.300000 l 78.000000 51.600000 l 78.300000 51.300000 l 79.200000 51.300000 l 79.500000 51.000000 l 80.400000 51.000000 l 81.000000 50.400000 l 81.300000 50.400000 l 81.600000 50.100000 l 81.900000 50.100000 l 82.200000 49.800000 l 82.500000 49.800000 l 83.700000 48.600000 l 83.700000 48.300000 l 84.000000 48.000000 l 84.000000 47.700000 l 84.300000 47.700000 l 84.600000 47.400000 l 84.300000 47.100000 l 84.600000 46.800000 l 84.300000 46.500000 l 84.300000 45.600000 l 84.000000 45.300000 l 84.000000 45.000000 l 83.400000 44.400000 l 83.400000 44.100000 l 82.800000 43.500000 l 82.200000 43.500000 l 81.600000 42.900000 l 81.300000 42.900000 l 81.000000 42.600000 l 80.700000 42.600000 l 80.400000 42.300000 l 79.800000 42.300000 l 79.500000 42.000000 l 78.600000 42.000000 l 78.300000 41.700000 l 77.700000 41.700000 l 77.700000 41.400000 l 77.400000 41.100000 l 77.100000 41.100000 l 76.800000 40.800000 l 75.900000 40.800000 l 75.600000 40.500000 l 75.000000 40.500000 l 74.700000 40.200000 l 74.400000 40.200000 l 74.100000 40.500000 l 73.800000 40.500000 l 73.500000 40.200000 l 71.400000 40.200000 l 71.100000 39.900000 l 70.200000 39.900000 l 69.900000 39.600000 l 66.000000 39.600000 l 65.700000 39.900000 l 65.400000 39.900000 l 65.400000 39.600000 l 65.100000 39.300000 l 57.900000 39.300000 l 57.300000 39.900000 l 54.000000 39.900000 l 53.700000 40.200000 l 53.400000 40.200000 l 53.100000 39.900000 l 51.600000 39.900000 l 51.300000 40.200000 l 50.700000 40.200000 l 50.400000 40.500000 l 49.200000 40.500000 l 48.600000 41.100000 l 47.400000 41.100000 l 47.100000 41.400000 l 46.200000 41.400000 l 45.900000 41.700000 l 45.300000 41.700000 l 44.700000 42.300000 l 43.800000 42.300000 l 43.500000 42.600000 l 43.200000 42.600000 l 42.900000 42.900000 l 42.600000 42.900000 l 42.300000 43.200000 l 42.300000 43.500000 l 42.000000 43.500000 l 41.700000 43.800000 l 41.400000 43.500000 l 39.900000 45.000000 l 39.600000 45.000000 l 39.300000 45.300000 l 39.300000 45.600000 l 39.000000 45.900000 l 38.700000 45.900000 l 38.400000 46.200000 l 39.000000 46.800000 l 39.000000 47.700000 l 39.600000 48.300000 l 39.900000 48.300000 l 40.800000 49.200000 l 41.100000 48.900000 l 41.100000 49.200000 l 42.000000 50.100000 l 42.300000 50.100000 l 42.600000 50.400000 l 42.900000 50.400000 l 43.500000 51.000000 l 44.100000 51.000000 l 44.400000 51.300000 l 45.300000 51.300000 l 45.600000 51.600000 l 45.900000 51.600000 l 46.200000 51.900000 l 46.800000 51.900000 l 47.100000 52.200000 l 47.700000 52.200000 l 48.000000 52.500000 l 48.900000 52.500000 l 49.200000 52.800000 l 50.700000 52.800000 l 51.000000 52.500000 l 51.300000 52.800000 l 51.600000 52.800000 l 51.900000 53.100000 l 52.500000 53.100000 l 52.800000 52.800000 l 53.100000 53.100000 l 53.700000 53.100000 l h 105.000000 101.400000 m 105.300000 101.700000 l 105.900000 101.100000 l 105.600000 100.800000 l 105.600000 100.500000 l 105.300000 100.200000 l 105.300000 99.900000 l 105.900000 99.300000 l 105.600000 99.000000 l 104.700000 99.000000 l 104.400000 99.300000 l 104.700000 99.600000 l 104.700000 100.800000 l 105.000000 101.100000 l h 105.000000 103.800000 m 105.300000 104.100000 l 105.600000 104.100000 l 105.900000 103.800000 l 105.900000 103.200000 l 105.600000 102.900000 l 105.300000 102.900000 l 105.000000 103.200000 l h 105.000000 109.200000 m 105.300000 109.500000 l 105.600000 109.200000 l 105.300000 108.900000 l h 105.300000 112.200000 m 105.600000 112.500000 l 106.200000 111.900000 l 106.200000 111.300000 l 105.900000 111.000000 l 105.900000 110.400000 l 106.200000 110.100000 l 105.900000 109.800000 l 105.900000 109.500000 l 105.600000 109.200000 l 105.000000 109.800000 l 105.000000 110.400000 l 105.300000 110.700000 l h 105.300000 116.100000 m 105.600000 116.400000 l 106.200000 115.800000 l 105.900000 115.500000 l 105.900000 115.200000 l 105.600000 114.900000 l 105.300000 115.200000 l h 105.000000 116.400000 m 105.300000 116.700000 l 105.600000 116.400000 l 105.300000 116.100000 l h 105.300000 120.900000 m 105.600000 121.200000 l 105.900000 120.900000 l 105.900000 120.600000 l 106.200000 120.300000 l 106.200000 119.400000 l 105.900000 119.100000 l 105.900000 118.800000 l 106.200000 118.500000 l 105.600000 117.900000 l 105.300000 118.200000 l h 105.900000 134.700000 m 106.200000 135.000000 l 106.500000 134.700000 l 106.500000 134.400000 l 105.900000 133.800000 l 105.600000 133.800000 l 105.300000 134.100000 l 105.300000 134.400000 l 105.600000 134.700000 l h 105.600000 135.000000 m 105.900000 135.300000 l 106.200000 135.000000 l 105.900000 134.700000 l h 104.400000 155.400000 m 105.000000 154.800000 l 105.000000 152.100000 l 105.300000 151.800000 l 105.300000 151.500000 l 105.000000 151.200000 l 105.000000 149.400000 l 104.700000 149.100000 l 104.400000 149.100000 l 104.100000 148.800000 l 102.300000 148.800000 l 102.000000 148.500000 l 101.700000 148.500000 l 101.700000 147.600000 l 102.000000 147.300000 l 102.300000 147.600000 l 103.200000 147.600000 l 103.500000 147.900000 l 103.800000 147.600000 l 104.100000 147.600000 l 105.000000 148.500000 l 105.600000 147.900000 l 105.300000 147.600000 l 105.000000 147.600000 l 105.300000 147.300000 l 104.700000 146.700000 l 105.000000 146.400000 l 105.000000 143.100000 l 104.700000 142.800000 l 104.700000 140.700000 l 105.000000 140.400000 l 105.300000 140.400000 l 105.600000 140.100000 l 105.300000 139.800000 l 105.000000 139.800000 l 105.000000 135.900000 l 104.700000 135.600000 l 104.700000 135.300000 l 105.000000 135.000000 l 104.400000 134.400000 l 104.700000 134.100000 l 104.700000 133.500000 l 104.400000 133.200000 l 104.400000 132.900000 l 104.700000 132.600000 l 104.400000 132.300000 l 104.700000 132.000000 l 104.700000 130.200000 l 104.400000 129.900000 l 104.700000 129.600000 l 104.400000 129.300000 l 104.400000 129.000000 l 105.000000 128.400000 l 105.000000 128.100000 l 104.700000 127.800000 l 104.700000 126.300000 l 104.400000 126.000000 l 104.700000 125.700000 l 104.700000 123.300000 l 104.400000 123.000000 l 104.700000 122.700000 l 104.400000 122.400000 l 104.400000 120.900000 l 104.700000 120.600000 l 104.700000 120.000000 l 104.400000 119.700000 l 104.700000 119.400000 l 104.400000 119.100000 l 104.700000 118.800000 l 104.700000 118.500000 l 104.400000 118.200000 l 104.400000 117.900000 l 103.800000 117.300000 l 104.100000 117.000000 l 104.100000 116.400000 l 104.400000 116.400000 l 104.700000 116.100000 l 104.400000 115.800000 l 104.400000 115.500000 l 104.700000 115.200000 l 104.700000 114.900000 l 104.400000 114.600000 l 104.400000 111.600000 l 104.700000 111.300000 l 104.400000 111.000000 l 104.400000 110.400000 l 104.700000 110.100000 l 104.400000 109.800000 l 104.400000 109.500000 l 104.100000 109.200000 l 104.100000 108.300000 l 104.400000 108.000000 l 104.100000 107.700000 l 104.100000 104.400000 l 104.400000 104.100000 l 104.400000 101.100000 l 104.100000 100.800000 l 104.100000 100.500000 l 103.800000 100.200000 l 104.100000 99.900000 l 104.100000 99.600000 l 104.400000 99.300000 l 104.100000 99.000000 l 103.200000 99.000000 l 102.900000 98.700000 l 102.600000 98.700000 l 102.300000 99.000000 l 100.800000 99.000000 l 100.500000 99.300000 l 99.900000 99.300000 l 99.900000 97.800000 l 100.500000 97.200000 l 100.500000 96.900000 l 100.800000 97.200000 l 101.400000 97.200000 l 101.700000 97.500000 l 102.000000 97.200000 l 102.300000 97.500000 l 102.600000 97.500000 l 103.200000 98.100000 l 103.800000 98.100000 l 104.100000 97.800000 l 103.800000 97.500000 l 103.800000 97.200000 l 104.100000 96.900000 l 103.800000 96.600000 l 103.800000 93.300000 l 104.100000 93.000000 l 103.500000 92.400000 l 103.200000 92.400000 l 103.500000 92.100000 l 103.500000 91.800000 l 103.800000 91.500000 l 103.800000 91.200000 l 103.500000 90.900000 l 103.500000 89.100000 l 103.800000 88.800000 l 103.500000 88.500000 l 103.800000 88.200000 l 103.500000 87.900000 l 103.500000 87.000000 l 103.800000 86.700000 l 103.800000 85.500000 l 103.500000 85.200000 l 103.800000 84.900000 l 103.800000 84.300000 l 103.200000 83.700000 l 103.500000 83.400000 l 103.500000 82.800000 l 103.200000 82.500000 l 103.500000 82.200000 l 103.500000 81.300000 l 103.200000 81.000000 l 103.200000 80.400000 l 102.900000 80.100000 l 103.800000 79.200000 l 103.800000 78.300000 l 103.200000 77.700000 l 103.500000 77.700000 l 103.800000 77.400000 l 103.500000 77.100000 l 103.500000 76.500000 l 103.800000 76.200000 l 103.800000 75.900000 l 103.200000 75.300000 l 103.200000 73.500000 l 103.500000 73.200000 l 103.500000 72.900000 l 103.200000 72.600000 l 103.200000 72.300000 l 103.500000 72.000000 l 103.500000 68.700000 l 103.200000 68.400000 l 103.200000 68.100000 l 103.500000 67.800000 l 103.200000 67.500000 l 103.200000 66.300000 l 103.500000 66.000000 l 103.500000 65.700000 l 103.200000 65.400000 l 103.500000 65.100000 l 103.500000 62.400000 l 103.200000 62.100000 l 103.200000 61.200000 l 103.500000 60.900000 l 103.500000 60.300000 l 103.200000 60.000000 l 103.500000 59.700000 l 102.900000 59.100000 l 102.900000 58.200000 l 103.200000 57.900000 l 103.200000 55.800000 l 103.500000 55.500000 l 103.500000 54.300000 l 103.200000 54.000000 l 103.500000 53.700000 l 103.500000 51.000000 l 103.200000 50.700000 l 103.200000 50.400000 l 103.500000 50.100000 l 103.200000 49.800000 l 103.200000 49.200000 l 103.500000 48.900000 l 103.500000 47.400000 l 103.200000 47.100000 l 102.000000 47.100000 l 101.700000 46.800000 l 99.600000 46.800000 l 99.300000 47.100000 l 98.400000 47.100000 l 98.100000 47.400000 l 97.500000 47.400000 l 97.200000 47.100000 l 96.900000 47.400000 l 96.600000 47.100000 l 96.300000 47.400000 l 95.100000 47.400000 l 94.800000 47.100000 l 94.500000 47.400000 l 91.500000 47.400000 l 91.200000 47.700000 l 90.000000 47.700000 l 89.700000 47.400000 l 85.800000 47.400000 l 85.500000 47.700000 l 85.500000 48.300000 l 85.200000 48.600000 l 85.200000 48.900000 l 84.600000 49.500000 l 84.600000 49.800000 l 84.300000 49.800000 l 83.700000 50.400000 l 83.400000 50.400000 l 82.800000 51.000000 l 82.800000 51.300000 l 82.200000 51.900000 l 81.600000 51.900000 l 81.300000 52.200000 l 81.000000 52.200000 l 80.400000 52.800000 l 79.500000 52.800000 l 79.200000 53.100000 l 78.600000 53.100000 l 78.300000 53.400000 l 78.000000 53.400000 l 77.700000 53.100000 l 77.400000 53.100000 l 77.100000 53.400000 l 76.500000 53.400000 l 76.200000 53.700000 l 74.700000 53.700000 l 74.400000 54.000000 l 74.400000 54.300000 l 72.900000 54.300000 l 72.600000 54.600000 l 70.200000 54.600000 l 69.900000 54.900000 l 67.800000 54.900000 l 67.500000 55.200000 l 66.300000 55.200000 l 66.000000 55.500000 l 65.700000 55.200000 l 57.900000 55.200000 l 57.600000 55.500000 l 56.100000 55.500000 l 55.800000 55.200000 l 53.700000 55.200000 l 53.400000 54.900000 l 52.500000 54.900000 l 52.200000 54.600000 l 51.000000 54.600000 l 50.700000 54.300000 l 49.800000 54.300000 l 49.500000 54.000000 l 49.200000 54.300000 l 48.300000 54.300000 l 48.000000 54.000000 l 47.100000 54.000000 l 46.800000 53.700000 l 46.200000 53.700000 l 45.900000 53.400000 l 45.300000 53.400000 l 45.000000 53.100000 l 44.400000 53.100000 l 43.800000 52.500000 l 43.200000 52.500000 l 42.600000 51.900000 l 42.000000 51.900000 l 41.700000 51.600000 l 41.400000 51.600000 l 41.100000 51.300000 l 40.500000 51.300000 l 39.600000 50.400000 l 39.300000 50.400000 l 39.000000 50.100000 l 39.000000 49.800000 l 38.700000 49.500000 l 38.400000 49.500000 l 38.400000 49.200000 l 38.100000 48.900000 l 38.100000 48.300000 l 38.400000 48.300000 l 38.700000 48.000000 l 38.400000 47.700000 l 35.700000 47.700000 l 35.400000 48.000000 l 35.100000 48.000000 l 34.800000 47.700000 l 34.500000 47.700000 l 34.200000 48.000000 l 33.600000 48.000000 l 33.300000 47.700000 l 33.000000 48.000000 l 32.700000 48.000000 l 32.400000 48.300000 l 31.800000 48.300000 l 31.500000 48.000000 l 31.200000 48.300000 l 29.700000 48.300000 l 29.100000 47.700000 l 28.800000 48.000000 l 24.300000 48.000000 l 24.000000 48.300000 l 23.700000 48.000000 l 23.400000 48.300000 l 23.400000 51.900000 l 23.700000 52.200000 l 23.700000 52.500000 l 24.000000 52.800000 l 23.700000 53.100000 l 23.700000 54.900000 l 23.400000 55.200000 l 23.700000 55.500000 l 23.700000 57.000000 l 24.000000 57.300000 l 23.400000 57.900000 l 23.400000 58.800000 l 23.700000 59.100000 l 23.400000 59.400000 l 23.400000 60.300000 l 23.700000 60.600000 l 23.700000 63.000000 l 23.400000 63.300000 l 23.700000 63.600000 l 23.400000 63.900000 l 23.700000 64.200000 l 23.700000 67.800000 l 24.000000 68.100000 l 24.000000 68.700000 l 24.300000 69.000000 l 24.300000 69.600000 l 24.000000 69.600000 l 23.700000 69.900000 l 23.700000 70.800000 l 23.400000 71.100000 l 24.000000 71.700000 l 24.000000 72.000000 l 23.700000 72.300000 l 24.000000 72.600000 l 24.000000 73.500000 l 23.700000 73.800000 l 23.700000 74.100000 l 24.000000 74.400000 l 24.000000 74.700000 l 23.400000 75.300000 l 23.700000 75.600000 l 23.700000 76.800000 l 23.400000 77.100000 l 23.700000 77.400000 l 24.000000 77.400000 l 24.000000 77.700000 l 23.700000 78.000000 l 24.000000 78.300000 l 24.000000 79.800000 l 23.700000 80.100000 l 24.000000 80.400000 l 24.000000 81.600000 l 23.700000 81.900000 l 23.700000 85.800000 l 24.000000 86.100000 l 24.000000 86.400000 l 23.700000 86.700000 l 23.700000 87.300000 l 24.000000 87.600000 l 24.300000 87.600000 l 24.600000 87.900000 l 25.200000 87.900000 l 26.400000 89.100000 l 26.700000 89.100000 l 27.300000 89.700000 l 27.300000 90.000000 l 28.500000 91.200000 l 28.800000 91.200000 l 29.100000 91.500000 l 29.400000 91.500000 l 29.700000 91.800000 l 29.700000 92.100000 l 30.300000 92.700000 l 30.300000 93.000000 l 30.600000 93.300000 l 30.900000 93.300000 l 31.200000 93.600000 l 31.200000 93.900000 l 31.500000 94.200000 l 31.800000 94.200000 l 32.400000 94.800000 l 32.700000 94.800000 l 33.300000 95.400000 l 33.900000 95.400000 l 35.700000 97.200000 l 35.700000 97.500000 l 36.000000 97.800000 l 37.200000 97.800000 l 37.800000 98.400000 l 37.800000 99.000000 l 38.100000 99.300000 l 38.700000 99.300000 l 38.700000 99.600000 l 39.000000 99.900000 l 39.300000 99.900000 l 39.300000 100.200000 l 40.800000 101.700000 l 41.100000 101.700000 l 41.400000 102.000000 l 42.000000 102.000000 l 42.300000 102.300000 l 42.600000 102.300000 l 44.400000 104.100000 l 44.700000 104.100000 l 45.000000 104.400000 l 45.300000 104.400000 l 45.900000 105.000000 l 45.900000 105.300000 l 45.600000 105.600000 l 45.900000 105.900000 l 46.500000 105.900000 l 47.400000 106.800000 l 47.700000 106.800000 l 48.000000 107.100000 l 48.300000 107.100000 l 48.300000 107.400000 l 48.900000 108.000000 l 48.600000 108.000000 l 48.300000 108.300000 l 48.600000 108.600000 l 49.500000 108.600000 l 49.200000 108.900000 l 49.500000 109.200000 l 49.800000 109.200000 l 50.100000 109.500000 l 50.400000 109.500000 l 50.700000 109.800000 l 51.000000 109.800000 l 51.000000 110.100000 l 51.300000 110.400000 l 51.900000 110.400000 l 51.900000 110.700000 l 52.200000 111.000000 l 52.500000 111.000000 l 53.100000 111.600000 l 53.400000 111.600000 l 53.400000 111.900000 l 54.900000 113.400000 l 54.900000 113.700000 l 55.500000 114.300000 l 55.800000 114.300000 l 56.700000 115.200000 l 57.000000 115.200000 l 57.000000 115.500000 l 57.600000 116.100000 l 57.900000 116.100000 l 58.500000 116.700000 l 58.800000 116.700000 l 59.100000 117.000000 l 59.400000 117.000000 l 59.400000 117.300000 l 60.300000 118.200000 l 60.600000 118.200000 l 60.600000 118.500000 l 60.900000 118.800000 l 61.500000 118.800000 l 62.400000 119.700000 l 62.700000 119.700000 l 62.700000 120.000000 l 63.300000 120.600000 l 63.600000 120.600000 l 65.700000 122.700000 l 66.000000 122.700000 l 67.800000 124.500000 l 68.100000 124.500000 l 69.000000 125.400000 l 69.300000 125.400000 l 70.500000 126.600000 l 70.800000 126.600000 l 71.100000 126.900000 l 71.400000 126.900000 l 73.200000 128.700000 l 73.500000 128.700000 l 73.500000 129.300000 l 74.400000 130.200000 l 74.700000 130.200000 l 75.000000 130.500000 l 75.000000 130.800000 l 75.300000 131.100000 l 75.600000 131.100000 l 75.900000 131.400000 l 76.200000 131.400000 l 76.500000 131.700000 l 76.800000 131.700000 l 77.400000 132.300000 l 77.700000 132.300000 l 77.700000 132.600000 l 78.600000 133.500000 l 78.900000 133.500000 l 79.200000 133.800000 l 79.500000 133.800000 l 79.200000 134.100000 l 79.800000 134.700000 l 79.800000 135.000000 l 80.100000 135.300000 l 80.700000 135.300000 l 81.000000 135.600000 l 81.300000 135.600000 l 81.600000 135.900000 l 81.600000 136.200000 l 81.900000 136.500000 l 82.500000 136.500000 l 83.100000 137.100000 l 83.100000 137.400000 l 83.700000 138.000000 l 84.000000 138.000000 l 84.600000 138.600000 l 84.900000 138.600000 l 85.200000 138.900000 l 85.500000 138.900000 l 86.400000 139.800000 l 86.400000 140.100000 l 86.700000 140.400000 l 87.000000 140.400000 l 87.300000 140.700000 l 87.300000 141.000000 l 88.500000 142.200000 l 88.800000 142.200000 l 89.100000 142.500000 l 89.100000 142.800000 l 89.400000 143.100000 l 89.700000 143.100000 l 90.000000 143.400000 l 90.300000 143.400000 l 90.900000 144.000000 l 91.200000 144.000000 l 91.500000 144.300000 l 91.800000 144.300000 l 91.800000 144.600000 l 92.100000 144.900000 l 92.100000 145.200000 l 92.700000 145.800000 l 93.300000 145.800000 l 94.500000 147.000000 l 94.800000 147.000000 l 94.800000 147.300000 l 95.400000 147.900000 l 95.700000 147.900000 l 96.600000 148.800000 l 96.900000 148.800000 l 97.200000 149.100000 l 97.200000 149.400000 l 97.500000 149.700000 l 97.800000 149.700000 l 97.800000 150.000000 l 98.100000 150.300000 l 98.400000 150.300000 l 99.000000 150.900000 l 99.300000 150.900000 l 99.300000 151.200000 l 99.600000 151.500000 l 99.900000 151.500000 l 100.200000 151.800000 l 100.500000 151.800000 l 102.300000 153.600000 l 102.600000 153.600000 l 102.600000 153.900000 l 102.900000 154.200000 l 103.200000 154.200000 l h 57.900000 160.800000 m 58.200000 161.100000 l 58.500000 160.800000 l 58.200000 160.500000 l h 24.600000 169.500000 m 24.900000 169.800000 l 25.200000 169.500000 l 25.500000 169.800000 l 25.800000 169.500000 l 26.100000 169.500000 l 26.400000 169.200000 l 26.700000 169.500000 l 27.600000 169.500000 l 27.900000 169.200000 l 31.800000 169.200000 l 32.100000 169.500000 l 33.300000 169.500000 l 33.600000 169.800000 l 33.900000 169.800000 l 34.200000 169.500000 l 34.500000 169.800000 l 34.800000 169.500000 l 34.800000 169.200000 l 38.400000 169.200000 l 39.000000 168.600000 l 39.000000 168.300000 l 39.300000 168.000000 l 39.300000 167.700000 l 39.600000 167.400000 l 39.600000 167.100000 l 40.200000 166.500000 l 40.500000 166.500000 l 40.800000 166.200000 l 40.800000 165.900000 l 41.100000 166.200000 l 42.000000 165.300000 l 42.300000 165.300000 l 43.200000 164.400000 l 43.500000 164.400000 l 43.800000 164.100000 l 44.700000 164.100000 l 45.000000 163.800000 l 45.600000 163.800000 l 45.900000 163.500000 l 46.500000 163.500000 l 46.800000 163.200000 l 47.100000 163.200000 l 47.400000 162.900000 l 48.300000 162.900000 l 48.600000 162.600000 l 49.200000 162.600000 l 49.500000 162.300000 l 49.800000 162.600000 l 50.100000 162.300000 l 50.400000 162.300000 l 50.700000 162.000000 l 51.300000 162.000000 l 51.600000 161.700000 l 51.900000 161.700000 l 52.200000 161.400000 l 52.800000 161.400000 l 53.100000 161.100000 l 53.400000 161.400000 l 53.700000 161.400000 l 54.000000 161.100000 l 55.200000 161.100000 l 55.500000 160.800000 l 55.800000 160.800000 l 56.400000 161.400000 l 56.700000 161.100000 l 57.000000 161.100000 l 57.300000 160.800000 l 57.300000 160.500000 l 57.300000 160.800000 l 57.600000 161.100000 l 58.200000 160.500000 l 58.500000 160.500000 l 58.500000 160.800000 l 58.800000 161.100000 l 59.100000 160.800000 l 59.400000 160.800000 l 59.700000 160.500000 l 63.000000 160.500000 l 63.300000 160.200000 l 63.600000 160.200000 l 63.900000 160.500000 l 64.500000 160.500000 l 64.800000 160.200000 l 65.100000 160.500000 l 68.700000 160.500000 l 69.000000 160.800000 l 69.300000 160.500000 l 69.600000 160.500000 l 69.900000 160.800000 l 71.400000 160.800000 l 71.700000 161.100000 l 72.000000 160.800000 l 72.300000 161.100000 l 72.600000 161.100000 l 72.900000 161.400000 l 74.400000 161.400000 l 74.700000 161.700000 l 75.000000 161.400000 l 75.600000 162.000000 l 75.900000 161.700000 l 76.200000 162.000000 l 76.800000 161.400000 l 77.100000 161.700000 l 77.400000 161.700000 l 77.700000 162.000000 l 78.000000 162.000000 l 78.300000 162.300000 l 79.500000 162.300000 l 79.800000 162.600000 l 80.100000 162.600000 l 80.700000 163.200000 l 81.300000 163.200000 l 81.600000 163.500000 l 81.900000 163.500000 l 82.200000 163.800000 l 82.800000 163.800000 l 83.400000 164.400000 l 84.000000 164.400000 l 84.300000 164.700000 l 84.600000 164.700000 l 84.900000 165.000000 l 85.500000 165.000000 l 85.500000 165.600000 l 86.100000 166.200000 l 86.400000 166.200000 l 86.700000 166.500000 l 86.700000 166.800000 l 87.000000 167.100000 l 87.300000 167.100000 l 87.600000 167.400000 l 87.600000 167.700000 l 87.900000 168.000000 l 87.900000 168.300000 l 88.200000 168.600000 l 88.500000 168.300000 l 88.800000 168.600000 l 89.100000 168.300000 l 89.700000 168.300000 l 90.000000 168.600000 l 90.900000 168.600000 l 91.200000 168.300000 l 92.400000 168.300000 l 93.000000 168.900000 l 93.600000 168.300000 l 95.100000 168.300000 l 95.400000 168.000000 l 95.700000 168.300000 l 99.600000 168.300000 l 99.900000 168.000000 l 100.200000 168.000000 l 100.500000 168.300000 l 100.800000 168.000000 l 101.100000 168.000000 l 101.400000 168.300000 l 101.700000 168.300000 l 102.000000 168.000000 l 102.900000 168.000000 l 103.200000 167.700000 l 103.500000 168.000000 l 103.800000 167.700000 l 104.100000 168.000000 l 104.700000 168.000000 l 105.000000 167.700000 l 105.000000 165.000000 l 105.300000 164.700000 l 105.300000 164.100000 l 105.000000 163.800000 l 105.300000 163.500000 l 105.000000 163.200000 l 105.300000 162.900000 l 105.300000 159.900000 l 105.000000 159.600000 l 105.000000 159.000000 l 104.700000 158.700000 l 104.400000 158.700000 l 104.400000 158.400000 l 102.900000 156.900000 l 102.600000 156.900000 l 102.600000 156.600000 l 102.000000 156.000000 l 101.400000 156.000000 l 100.200000 154.800000 l 99.900000 154.800000 l 98.100000 153.000000 l 97.800000 153.000000 l 97.800000 152.700000 l 96.600000 151.500000 l 96.300000 151.500000 l 96.000000 151.200000 l 95.700000 151.200000 l 94.200000 149.700000 l 93.900000 149.700000 l 93.000000 148.800000 l 92.700000 148.800000 l 92.100000 148.200000 l 92.100000 147.900000 l 91.800000 147.600000 l 91.500000 147.600000 l 89.400000 145.500000 l 89.100000 145.500000 l 87.900000 144.300000 l 87.600000 144.300000 l 86.400000 143.100000 l 86.100000 143.100000 l 84.600000 141.600000 l 84.300000 141.900000 l 83.700000 141.300000 l 83.700000 141.000000 l 82.200000 139.500000 l 81.900000 139.500000 l 80.100000 137.700000 l 80.100000 137.400000 l 79.800000 137.100000 l 79.500000 137.100000 l 79.500000 136.800000 l 79.200000 136.500000 l 78.900000 136.500000 l 77.400000 135.000000 l 76.800000 135.000000 l 76.200000 134.400000 l 75.900000 134.400000 l 74.400000 132.900000 l 74.100000 132.900000 l 72.000000 130.800000 l 71.700000 130.800000 l 69.000000 128.100000 l 68.700000 128.100000 l 68.700000 127.800000 l 67.500000 126.600000 l 66.900000 126.600000 l 66.300000 126.000000 l 66.000000 126.000000 l 66.000000 125.700000 l 65.700000 125.400000 l 65.400000 125.400000 l 65.400000 125.100000 l 65.100000 124.800000 l 64.800000 124.800000 l 64.200000 124.200000 l 63.900000 124.200000 l 63.900000 123.900000 l 63.600000 123.600000 l 63.300000 123.600000 l 60.600000 120.900000 l 60.300000 120.900000 l 59.100000 119.700000 l 58.800000 119.700000 l 58.200000 119.100000 l 57.900000 119.100000 l 57.900000 118.800000 l 57.600000 118.500000 l 57.000000 118.500000 l 56.400000 117.900000 l 56.100000 117.900000 l 55.800000 117.600000 l 55.800000 117.300000 l 55.200000 116.700000 l 55.200000 116.400000 l 54.000000 115.200000 l 53.700000 115.200000 l 53.700000 114.900000 l 52.800000 114.000000 l 52.500000 114.000000 l 52.200000 113.700000 l 51.900000 113.700000 l 50.700000 112.500000 l 50.100000 112.500000 l 48.600000 111.000000 l 48.300000 111.000000 l 47.700000 110.400000 l 47.400000 110.400000 l 47.100000 110.100000 l 46.800000 110.100000 l 45.300000 108.600000 l 45.000000 108.600000 l 43.500000 107.100000 l 43.200000 107.100000 l 42.900000 106.800000 l 42.900000 106.500000 l 41.700000 105.300000 l 41.400000 105.300000 l 38.400000 102.300000 l 38.100000 102.300000 l 37.800000 102.000000 l 37.500000 102.000000 l 36.300000 100.800000 l 36.000000 100.800000 l 36.000000 100.500000 l 35.400000 99.900000 l 35.100000 99.900000 l 33.600000 98.400000 l 33.300000 98.700000 l 31.500000 96.900000 l 31.200000 96.900000 l 31.200000 96.600000 l 30.600000 96.000000 l 30.600000 95.700000 l 30.000000 95.100000 l 29.700000 95.100000 l 28.500000 93.900000 l 28.200000 93.900000 l 28.200000 93.600000 l 27.600000 93.000000 l 27.300000 93.000000 l 24.900000 90.600000 l 24.000000 90.600000 l 23.700000 90.900000 l 23.700000 91.200000 l 24.000000 91.500000 l 23.700000 91.800000 l 23.700000 92.100000 l 24.000000 92.400000 l 24.000000 93.300000 l 23.700000 93.600000 l 24.000000 93.900000 l 24.000000 96.300000 l 23.700000 96.600000 l 24.000000 96.900000 l 24.000000 97.500000 l 23.700000 97.800000 l 24.000000 98.100000 l 24.000000 100.800000 l 24.300000 101.100000 l 23.700000 101.700000 l 23.700000 102.000000 l 24.300000 102.600000 l 24.300000 103.200000 l 24.000000 103.500000 l 24.300000 103.800000 l 24.300000 106.200000 l 23.700000 106.800000 l 24.000000 107.100000 l 24.000000 109.200000 l 23.700000 109.500000 l 23.700000 110.100000 l 24.000000 110.400000 l 24.000000 110.700000 l 24.300000 111.000000 l 24.300000 111.300000 l 24.000000 111.600000 l 24.300000 111.900000 l 24.300000 112.500000 l 24.000000 112.800000 l 24.300000 113.100000 l 24.300000 114.300000 l 24.600000 114.600000 l 24.600000 115.500000 l 24.300000 115.800000 l 24.300000 116.100000 l 24.000000 116.400000 l 24.300000 116.700000 l 24.300000 118.200000 l 24.000000 118.500000 l 24.000000 119.100000 l 24.300000 119.400000 l 24.300000 120.600000 l 24.000000 120.900000 l 24.300000 121.200000 l 24.300000 121.800000 l 24.600000 122.100000 l 24.300000 122.400000 l 24.300000 123.300000 l 24.000000 123.600000 l 24.000000 125.700000 l 24.300000 126.000000 l 24.000000 126.300000 l 24.000000 126.900000 l 24.300000 127.200000 l 24.000000 127.500000 l 24.300000 127.800000 l 24.300000 128.700000 l 24.600000 129.000000 l 24.600000 129.600000 l 24.300000 129.900000 l 24.300000 130.200000 l 24.000000 130.500000 l 24.000000 131.400000 l 24.300000 131.700000 l 24.300000 132.000000 l 24.000000 132.300000 l 24.000000 133.500000 l 24.300000 133.800000 l 24.300000 134.100000 l 24.000000 134.400000 l 24.000000 134.700000 l 24.300000 135.000000 l 24.000000 135.300000 l 24.000000 135.900000 l 24.300000 136.200000 l 24.300000 136.800000 l 24.600000 137.100000 l 24.300000 137.400000 l 24.300000 142.500000 l 24.000000 142.800000 l 24.000000 144.000000 l 24.300000 144.300000 l 24.300000 145.500000 l 24.600000 145.800000 l 24.300000 146.100000 l 24.300000 147.900000 l 24.600000 148.200000 l 24.300000 148.500000 l 24.300000 150.000000 l 24.600000 150.300000 l 24.600000 150.600000 l 24.300000 150.900000 l 24.300000 151.200000 l 24.600000 151.500000 l 24.600000 151.800000 l 24.900000 152.100000 l 24.600000 152.100000 l 24.300000 152.400000 l 24.300000 152.700000 l 24.600000 153.000000 l 24.600000 153.600000 l 24.900000 153.900000 l 24.900000 154.200000 l 24.000000 155.100000 l 24.600000 155.700000 l 24.600000 156.000000 l 24.300000 156.300000 l 24.300000 156.600000 l 24.600000 156.900000 l 24.300000 157.200000 l 24.300000 158.400000 l 24.600000 158.700000 l 24.600000 159.000000 l 24.300000 159.300000 l 24.300000 159.600000 l 24.600000 159.900000 l 24.600000 160.500000 l 24.300000 160.800000 l 24.300000 161.100000 l 24.900000 161.700000 l 24.600000 162.000000 l 24.600000 162.300000 l 24.900000 162.600000 l 24.900000 166.500000 l 24.600000 166.500000 l 24.300000 166.800000 l 24.600000 167.100000 l 24.600000 167.400000 l 24.300000 167.700000 l 24.600000 168.000000 l h 58.200000 177.000000 m 58.500000 177.300000 l 58.800000 177.000000 l 59.100000 177.300000 l 59.700000 177.300000 l 60.000000 177.000000 l 60.600000 177.000000 l 60.900000 176.700000 l 61.200000 177.000000 l 61.800000 177.000000 l 62.100000 177.300000 l 62.400000 177.000000 l 63.300000 177.000000 l 63.600000 177.300000 l 63.900000 177.000000 l 64.200000 177.000000 l 64.500000 177.300000 l 64.800000 177.000000 l 66.000000 177.000000 l 66.300000 176.700000 l 66.600000 177.000000 l 66.900000 176.700000 l 67.200000 177.000000 l 67.500000 177.000000 l 67.800000 176.700000 l 68.100000 176.700000 l 68.400000 176.400000 l 68.700000 176.700000 l 69.900000 176.700000 l 70.200000 176.400000 l 71.100000 176.400000 l 71.400000 176.700000 l 71.700000 176.400000 l 72.000000 176.400000 l 72.300000 176.100000 l 73.800000 176.100000 l 74.100000 175.800000 l 75.300000 175.800000 l 75.600000 175.500000 l 76.200000 175.500000 l 76.500000 175.800000 l 77.400000 174.900000 l 78.300000 174.900000 l 78.900000 174.300000 l 80.100000 174.300000 l 80.400000 174.000000 l 80.700000 174.300000 l 81.600000 174.300000 l 81.900000 174.000000 l 81.900000 173.700000 l 82.200000 174.000000 l 82.500000 173.700000 l 82.800000 173.700000 l 83.100000 173.400000 l 83.400000 173.400000 l 83.700000 173.100000 l 84.000000 173.100000 l 85.200000 171.900000 l 85.200000 171.600000 l 85.500000 171.300000 l 85.800000 171.300000 l 86.700000 170.400000 l 86.700000 169.800000 l 86.400000 169.500000 l 86.400000 168.600000 l 86.100000 168.300000 l 85.800000 168.300000 l 85.200000 167.700000 l 84.900000 167.700000 l 84.600000 167.400000 l 84.300000 167.400000 l 83.400000 166.500000 l 83.100000 166.500000 l 82.800000 166.200000 l 82.500000 166.200000 l 82.200000 165.900000 l 81.900000 165.900000 l 81.600000 165.600000 l 81.000000 165.600000 l 80.700000 165.300000 l 80.400000 165.300000 l 79.500000 164.400000 l 78.900000 164.400000 l 78.600000 164.100000 l 77.700000 164.100000 l 77.400000 163.800000 l 77.100000 163.800000 l 76.800000 163.500000 l 76.500000 163.800000 l 75.600000 163.800000 l 75.300000 163.500000 l 74.700000 163.500000 l 74.400000 163.200000 l 72.300000 163.200000 l 72.000000 162.900000 l 69.900000 162.900000 l 69.600000 162.600000 l 65.700000 162.600000 l 65.400000 162.300000 l 63.000000 162.300000 l 62.700000 162.600000 l 60.000000 162.600000 l 59.700000 162.900000 l 58.800000 162.900000 l 58.500000 163.200000 l 56.400000 163.200000 l 56.100000 162.900000 l 55.200000 162.900000 l 54.900000 163.200000 l 53.100000 163.200000 l 52.800000 163.500000 l 51.900000 163.500000 l 51.300000 164.100000 l 50.400000 164.100000 l 50.100000 164.400000 l 49.200000 164.400000 l 48.900000 164.700000 l 48.300000 164.700000 l 48.000000 165.000000 l 47.700000 164.700000 l 47.400000 165.000000 l 47.100000 165.000000 l 46.800000 165.300000 l 46.500000 165.300000 l 46.200000 165.600000 l 45.600000 165.600000 l 45.300000 165.900000 l 44.700000 165.900000 l 44.400000 166.200000 l 43.200000 166.200000 l 42.900000 166.500000 l 42.900000 166.800000 l 42.600000 167.100000 l 42.300000 167.100000 l 41.700000 167.700000 l 41.400000 167.700000 l 41.100000 168.000000 l 41.100000 168.300000 l 40.500000 168.900000 l 40.500000 170.400000 l 40.800000 170.700000 l 40.800000 171.300000 l 41.100000 171.600000 l 41.400000 171.600000 l 41.400000 171.900000 l 41.700000 172.200000 l 42.000000 172.200000 l 42.600000 172.800000 l 43.500000 172.800000 l 43.800000 173.100000 l 44.100000 173.100000 l 44.100000 173.400000 l 44.400000 173.700000 l 44.700000 173.700000 l 45.000000 174.000000 l 45.300000 174.000000 l 45.600000 174.300000 l 46.200000 174.300000 l 46.500000 174.600000 l 47.100000 174.600000 l 47.400000 174.900000 l 48.000000 174.900000 l 48.300000 175.200000 l 48.600000 175.200000 l 48.900000 175.500000 l 49.200000 175.200000 l 49.500000 175.500000 l 49.800000 175.500000 l 50.100000 175.800000 l 51.000000 175.800000 l 51.300000 176.100000 l 51.900000 175.500000 l 52.200000 175.800000 l 52.500000 175.800000 l 52.800000 176.100000 l 53.100000 175.800000 l 53.400000 176.100000 l 54.600000 176.100000 l 54.900000 176.400000 l 55.200000 176.100000 l 55.500000 176.400000 l 56.100000 176.400000 l 56.400000 176.700000 l 56.700000 176.700000 l 57.000000 177.000000 l h f* Q}\special{pdf:econtent}}
\put(0,0){\special{pdf:bcontent}\special{pdf:code q 0 g 184.200000 118.200000 m 183.900000 117.900000 l 183.600000 117.900000 l 183.300000 117.600000 l 183.000000 117.600000 l 183.000000 115.800000 l 182.700000 115.500000 l 183.000000 115.200000 l 183.000000 114.900000 l 183.300000 114.600000 l 183.000000 114.300000 l 183.300000 114.000000 l 183.300000 113.700000 l 183.000000 113.400000 l 183.300000 113.100000 l 183.300000 111.300000 l 183.000000 111.000000 l 183.000000 110.700000 l 182.700000 110.400000 l 180.600000 110.400000 l 180.300000 110.100000 l 178.500000 110.100000 l 178.200000 110.400000 l 176.700000 110.400000 l 176.400000 110.100000 l 176.100000 110.400000 l 176.100000 110.700000 l 176.100000 110.400000 l 175.800000 110.100000 l 175.500000 110.100000 l 175.200000 110.400000 l 172.500000 110.400000 l 172.200000 110.100000 l 171.900000 110.400000 l 169.200000 110.400000 l 168.900000 110.100000 l 168.000000 110.100000 l 167.700000 110.400000 l 167.400000 110.100000 l 167.100000 110.100000 l 166.800000 109.800000 l 166.800000 109.200000 l 166.500000 108.900000 l 166.500000 108.600000 l 165.900000 108.000000 l 165.600000 108.000000 l 165.600000 106.800000 l 165.900000 106.500000 l 165.600000 106.200000 l 165.600000 105.900000 l 165.900000 105.600000 l 165.900000 104.400000 l 165.600000 104.100000 l 165.600000 103.200000 l 165.900000 103.200000 l 166.200000 102.900000 l 166.200000 102.600000 l 166.500000 102.300000 l 166.500000 102.000000 l 166.200000 101.700000 l 166.200000 100.200000 l 165.900000 99.900000 l 165.900000 98.400000 l 166.200000 98.100000 l 165.900000 97.800000 l 165.900000 96.600000 l 165.600000 96.300000 l 165.900000 96.000000 l 165.900000 93.900000 l 165.600000 93.600000 l 166.200000 93.000000 l 166.200000 92.700000 l 165.900000 92.400000 l 165.900000 91.800000 l 165.600000 91.500000 l 165.600000 90.900000 l 165.900000 90.600000 l 165.900000 84.600000 l 166.200000 84.300000 l 166.200000 84.000000 l 166.800000 83.400000 l 166.500000 83.100000 l 166.500000 82.800000 l 166.800000 82.800000 l 167.700000 81.900000 l 168.000000 81.900000 l 168.300000 81.600000 l 170.700000 81.600000 l 171.000000 81.300000 l 173.100000 81.300000 l 173.400000 81.000000 l 173.700000 81.300000 l 174.600000 81.300000 l 174.900000 81.000000 l 175.500000 81.600000 l 176.100000 81.000000 l 176.400000 81.300000 l 179.700000 81.300000 l 180.000000 81.000000 l 180.300000 81.300000 l 180.600000 81.300000 l 180.900000 81.000000 l 181.200000 81.300000 l 181.800000 81.300000 l 182.100000 81.000000 l 182.700000 81.000000 l 183.000000 81.300000 l 183.300000 81.000000 l 183.600000 81.000000 l 183.900000 81.300000 l 184.500000 80.700000 l 186.300000 80.700000 l 186.600000 81.000000 l 186.900000 80.700000 l 187.200000 80.700000 l 187.500000 81.000000 l 188.400000 81.000000 l 188.700000 81.300000 l 189.000000 81.000000 l 192.300000 81.000000 l 192.600000 80.700000 l 193.800000 80.700000 l 194.100000 80.400000 l 194.400000 80.700000 l 197.400000 80.700000 l 197.700000 80.400000 l 198.000000 80.700000 l 199.200000 80.700000 l 199.500000 81.000000 l 200.700000 81.000000 l 201.000000 80.700000 l 201.600000 80.700000 l 201.900000 80.400000 l 202.500000 80.400000 l 202.800000 80.700000 l 204.000000 80.700000 l 204.300000 80.400000 l 204.900000 80.400000 l 205.200000 80.100000 l 205.500000 80.400000 l 209.100000 80.400000 l 209.400000 80.100000 l 210.000000 80.100000 l 210.300000 79.800000 l 210.600000 80.100000 l 211.200000 80.100000 l 211.800000 80.700000 l 212.700000 80.700000 l 213.000000 81.000000 l 213.300000 81.000000 l 213.900000 81.600000 l 214.200000 81.600000 l 215.100000 82.500000 l 215.400000 82.500000 l 215.400000 82.800000 l 215.100000 83.100000 l 215.700000 83.700000 l 215.700000 84.900000 l 216.000000 85.200000 l 216.000000 85.500000 l 215.700000 85.800000 l 216.000000 86.100000 l 216.000000 86.400000 l 216.300000 86.700000 l 216.000000 87.000000 l 216.000000 87.300000 l 215.700000 87.600000 l 216.000000 87.900000 l 215.700000 88.200000 l 215.700000 88.500000 l 216.300000 89.100000 l 215.700000 89.700000 l 215.700000 90.600000 l 216.000000 90.900000 l 216.000000 92.400000 l 216.300000 92.700000 l 216.000000 93.000000 l 216.300000 93.300000 l 216.300000 93.900000 l 216.000000 94.200000 l 216.300000 94.500000 l 216.000000 94.800000 l 216.300000 95.100000 l 216.300000 97.800000 l 216.000000 98.100000 l 216.000000 98.400000 l 216.300000 98.700000 l 216.300000 99.600000 l 216.600000 99.900000 l 216.300000 100.200000 l 216.600000 100.500000 l 216.600000 101.700000 l 216.300000 102.000000 l 216.300000 102.600000 l 216.000000 102.900000 l 216.000000 103.500000 l 216.300000 103.800000 l 216.300000 104.100000 l 216.000000 104.400000 l 216.300000 104.700000 l 216.300000 105.000000 l 216.000000 105.300000 l 216.300000 105.600000 l 216.000000 105.600000 l 215.700000 105.900000 l 216.300000 106.500000 l 216.300000 106.800000 l 215.700000 107.400000 l 215.700000 108.000000 l 215.100000 108.600000 l 215.400000 108.900000 l 215.100000 109.200000 l 215.100000 109.500000 l 215.100000 109.200000 l 214.800000 108.900000 l 214.500000 108.900000 l 214.200000 109.200000 l 213.000000 109.200000 l 212.700000 109.500000 l 212.400000 109.500000 l 212.100000 109.200000 l 210.000000 109.200000 l 209.700000 109.500000 l 203.700000 109.500000 l 203.400000 109.800000 l 202.500000 109.800000 l 202.200000 109.500000 l 201.600000 109.500000 l 199.800000 111.300000 l 199.800000 111.900000 l 200.100000 112.200000 l 199.800000 112.500000 l 200.100000 112.800000 l 200.100000 113.100000 l 199.800000 113.400000 l 199.800000 114.000000 l 199.500000 114.300000 l 199.800000 114.600000 l 199.800000 115.800000 l 199.500000 116.100000 l 199.800000 116.400000 l 199.800000 116.700000 l 200.100000 117.000000 l 200.100000 117.300000 l 199.500000 117.300000 l 199.200000 117.600000 l 198.600000 117.600000 l 198.300000 117.900000 l 198.000000 117.600000 l 196.200000 117.600000 l 195.900000 117.900000 l 195.600000 117.600000 l 193.200000 117.600000 l 192.600000 118.200000 l 188.400000 118.200000 l 188.100000 117.900000 l 184.800000 117.900000 l 184.500000 118.200000 l h 209.100000 82.800000 m 209.400000 83.100000 l 211.800000 83.100000 l 212.100000 82.800000 l 211.800000 82.500000 l 209.400000 82.500000 l h 199.500000 82.800000 m 199.800000 83.100000 l 200.700000 83.100000 l 201.000000 82.800000 l 200.700000 82.500000 l 199.800000 82.500000 l h 211.800000 83.100000 m 212.100000 83.400000 l 212.700000 83.400000 l 213.000000 83.100000 l 212.700000 82.800000 l 212.100000 82.800000 l h 201.300000 83.700000 m 202.200000 83.700000 l 202.500000 83.400000 l 204.900000 83.400000 l 205.200000 83.700000 l 205.500000 83.400000 l 207.000000 83.400000 l 207.300000 83.100000 l 207.000000 82.800000 l 204.000000 82.800000 l 204.000000 82.500000 l 203.700000 82.200000 l 201.000000 82.200000 l 200.700000 82.500000 l 201.000000 82.800000 l 202.500000 82.800000 l 201.000000 82.800000 l 200.700000 83.100000 l h 199.800000 83.700000 m 200.100000 84.000000 l 200.700000 84.000000 l 201.000000 83.700000 l 201.000000 83.400000 l 200.700000 83.100000 l 199.500000 83.100000 l 199.200000 82.800000 l 198.600000 82.800000 l 198.300000 83.100000 l 198.000000 82.800000 l 197.400000 83.400000 l 197.700000 83.700000 l h 168.000000 84.600000 m 168.300000 84.300000 l 183.600000 84.300000 l 183.900000 84.000000 l 184.200000 84.000000 l 184.500000 83.700000 l 185.700000 83.700000 l 186.000000 84.000000 l 186.600000 84.000000 l 186.900000 83.700000 l 187.500000 83.700000 l 188.100000 84.300000 l 188.400000 84.000000 l 192.300000 84.000000 l 192.600000 83.700000 l 197.100000 83.700000 l 197.400000 83.400000 l 197.100000 83.100000 l 196.500000 83.100000 l 196.500000 82.800000 l 198.600000 82.800000 l 198.900000 82.500000 l 198.600000 82.200000 l 197.700000 82.200000 l 197.400000 81.900000 l 197.100000 82.200000 l 193.500000 82.200000 l 193.200000 82.500000 l 192.900000 82.200000 l 192.600000 82.200000 l 192.300000 82.500000 l 189.000000 82.500000 l 188.700000 82.800000 l 188.400000 82.800000 l 188.100000 82.500000 l 187.500000 82.500000 l 187.200000 82.200000 l 185.100000 82.200000 l 184.800000 82.500000 l 184.500000 82.200000 l 184.200000 82.200000 l 183.600000 82.800000 l 182.100000 82.800000 l 181.800000 82.500000 l 181.500000 82.800000 l 181.200000 82.800000 l 180.900000 82.500000 l 180.600000 82.800000 l 179.400000 82.800000 l 179.100000 82.500000 l 178.800000 82.800000 l 178.500000 82.800000 l 178.200000 82.500000 l 177.900000 82.500000 l 177.600000 82.800000 l 177.300000 82.500000 l 177.000000 82.800000 l 168.900000 82.800000 l 168.600000 83.100000 l 168.300000 83.100000 l 167.400000 84.000000 l h 198.000000 87.600000 m 198.300000 87.900000 l 199.500000 87.900000 l 199.800000 87.600000 l 199.500000 87.300000 l 198.300000 87.300000 l h 187.200000 87.600000 m 187.500000 87.900000 l 187.800000 87.600000 l 187.500000 87.300000 l h 211.800000 88.200000 m 212.100000 88.500000 l 213.600000 88.500000 l 213.900000 88.200000 l 213.600000 87.900000 l 212.100000 87.900000 l h 197.700000 88.200000 m 198.000000 88.500000 l 198.300000 88.200000 l 198.000000 87.900000 l h 180.600000 88.500000 m 180.900000 88.200000 l 180.900000 87.900000 l 180.600000 87.600000 l 180.300000 87.600000 l 180.000000 87.900000 l h 176.700000 88.200000 m 177.000000 88.500000 l 177.300000 88.200000 l 177.000000 87.900000 l h 203.400000 88.500000 m 203.700000 88.800000 l 204.000000 88.500000 l 203.700000 88.200000 l h 197.400000 88.500000 m 197.700000 88.800000 l 198.000000 88.500000 l 197.700000 88.200000 l h 207.300000 88.800000 m 207.600000 89.100000 l 207.900000 88.800000 l 207.600000 88.500000 l h 196.200000 89.100000 m 196.500000 89.400000 l 196.800000 89.100000 l 196.500000 88.800000 l h 210.300000 89.400000 m 210.600000 89.700000 l 211.200000 89.100000 l 210.900000 88.800000 l 209.400000 88.800000 l 209.100000 89.100000 l 209.400000 89.400000 l h 205.200000 89.400000 m 205.500000 89.700000 l 206.700000 89.700000 l 207.000000 89.400000 l 207.300000 89.400000 l 207.600000 89.100000 l 207.300000 88.800000 l 206.700000 88.800000 l 206.400000 89.100000 l 206.400000 88.800000 l 206.100000 88.500000 l 205.800000 88.500000 l 205.500000 88.800000 l 205.500000 89.100000 l h 203.700000 89.400000 m 204.000000 89.700000 l 205.500000 88.200000 l 205.200000 87.900000 l 204.900000 87.900000 l 203.700000 89.100000 l h 201.300000 89.400000 m 201.600000 89.700000 l 203.400000 89.700000 l 203.700000 89.400000 l 203.700000 88.800000 l 203.400000 88.500000 l 202.800000 89.100000 l 202.200000 89.100000 l 201.900000 88.800000 l h 197.700000 89.400000 m 198.000000 89.700000 l 199.200000 89.700000 l 199.500000 89.400000 l 200.700000 89.400000 l 201.000000 89.100000 l 201.000000 88.800000 l 201.900000 87.900000 l 205.200000 87.900000 l 205.500000 88.200000 l 205.800000 87.900000 l 206.100000 87.900000 l 206.400000 87.600000 l 207.600000 87.600000 l 207.900000 87.900000 l 208.200000 87.900000 l 208.500000 87.600000 l 208.800000 87.600000 l 209.100000 87.300000 l 208.800000 87.000000 l 207.900000 87.000000 l 207.600000 86.700000 l 207.300000 86.700000 l 207.000000 87.000000 l 206.100000 87.000000 l 205.800000 86.700000 l 205.500000 87.000000 l 204.000000 87.000000 l 203.700000 87.300000 l 202.500000 87.300000 l 202.500000 87.000000 l 202.200000 86.700000 l 201.900000 87.000000 l 201.000000 87.000000 l 200.400000 87.600000 l 199.800000 87.600000 l 199.500000 87.900000 l 199.500000 88.200000 l 199.800000 88.500000 l 199.500000 88.500000 l 199.200000 88.800000 l 198.600000 88.800000 l 198.300000 89.100000 l 198.000000 89.100000 l h 195.900000 89.400000 m 196.200000 89.700000 l 196.500000 89.400000 l 196.200000 89.100000 l h 196.500000 89.700000 m 196.800000 90.000000 l 197.100000 89.700000 l 197.400000 89.700000 l 197.700000 89.400000 l 197.700000 88.800000 l 197.400000 88.500000 l 197.100000 88.500000 l 196.800000 88.800000 l 196.800000 89.100000 l 196.500000 89.400000 l h 192.900000 90.000000 m 193.200000 90.300000 l 193.500000 90.000000 l 193.200000 89.700000 l h 168.000000 90.300000 m 168.300000 90.000000 l 171.000000 90.000000 l 171.300000 90.300000 l 171.600000 90.000000 l 172.800000 90.000000 l 173.100000 89.700000 l 173.700000 90.300000 l 174.000000 90.300000 l 174.300000 90.000000 l 174.900000 90.000000 l 175.500000 89.400000 l 175.500000 88.800000 l 175.200000 88.500000 l 174.600000 88.500000 l 174.300000 88.800000 l 173.100000 88.800000 l 172.800000 88.500000 l 172.800000 88.200000 l 173.400000 87.600000 l 173.100000 87.300000 l 167.700000 87.300000 l 167.400000 87.600000 l 167.700000 87.300000 l 167.700000 87.000000 l 173.400000 87.000000 l 173.400000 87.300000 l 173.700000 87.600000 l 174.000000 87.300000 l 174.000000 87.000000 l 174.300000 86.700000 l 175.500000 86.700000 l 176.100000 86.100000 l 180.600000 86.100000 l 180.600000 86.400000 l 180.900000 86.700000 l 181.200000 86.700000 l 181.500000 86.400000 l 181.800000 86.400000 l 182.100000 86.100000 l 183.900000 86.100000 l 184.200000 85.800000 l 184.500000 86.100000 l 187.500000 86.100000 l 187.800000 86.400000 l 187.500000 86.100000 l 184.500000 86.100000 l 184.200000 85.800000 l 183.600000 86.400000 l 181.800000 86.400000 l 181.500000 86.700000 l 181.800000 87.000000 l 182.100000 87.000000 l 179.400000 87.000000 l 179.100000 86.700000 l 178.800000 87.000000 l 178.500000 87.000000 l 178.200000 86.700000 l 177.300000 86.700000 l 177.000000 87.000000 l 176.100000 87.000000 l 175.800000 87.300000 l 175.500000 87.300000 l 175.200000 87.000000 l 174.900000 87.300000 l 174.300000 87.300000 l 174.000000 87.600000 l 174.000000 87.900000 l 174.300000 88.200000 l 175.200000 88.200000 l 175.500000 88.500000 l 175.800000 88.500000 l 176.100000 88.200000 l 176.700000 88.200000 l 177.000000 87.900000 l 177.900000 87.900000 l 178.200000 88.200000 l 178.800000 87.600000 l 179.100000 87.900000 l 179.400000 87.600000 l 179.700000 87.600000 l 180.000000 87.900000 l 180.300000 87.600000 l 180.600000 87.600000 l 181.200000 88.200000 l 181.500000 87.900000 l 181.500000 87.600000 l 183.000000 87.600000 l 183.300000 87.300000 l 183.900000 87.300000 l 184.200000 87.000000 l 186.000000 87.000000 l 186.300000 87.300000 l 187.500000 87.300000 l 187.800000 87.600000 l 188.100000 87.300000 l 189.600000 87.300000 l 189.900000 87.600000 l 187.800000 87.600000 l 187.500000 87.900000 l 187.500000 88.200000 l 187.800000 88.500000 l 190.800000 88.500000 l 191.100000 88.200000 l 191.700000 88.200000 l 192.000000 88.500000 l 192.300000 88.500000 l 192.600000 88.200000 l 193.200000 88.200000 l 193.500000 87.900000 l 195.300000 87.900000 l 195.600000 87.600000 l 196.200000 87.600000 l 196.800000 87.000000 l 198.600000 87.000000 l 199.200000 86.400000 l 198.900000 86.100000 l 192.600000 86.100000 l 192.300000 86.400000 l 188.700000 86.400000 l 190.200000 86.400000 l 190.500000 86.100000 l 190.800000 86.400000 l 191.100000 86.400000 l 191.400000 86.100000 l 191.700000 86.400000 l 192.300000 86.400000 l 192.600000 86.100000 l 197.100000 86.100000 l 197.400000 85.800000 l 197.700000 86.100000 l 198.000000 85.800000 l 198.300000 86.100000 l 198.600000 85.800000 l 199.500000 85.800000 l 199.800000 86.100000 l 200.700000 86.100000 l 201.000000 85.800000 l 201.300000 86.100000 l 201.900000 86.100000 l 202.200000 85.800000 l 208.500000 85.800000 l 209.100000 86.400000 l 209.700000 86.400000 l 210.000000 86.700000 l 210.600000 86.700000 l 209.400000 86.700000 l 209.100000 87.000000 l 209.400000 87.300000 l 211.500000 87.300000 l 211.800000 87.600000 l 212.100000 87.300000 l 213.900000 87.300000 l 214.200000 87.000000 l 214.200000 86.400000 l 213.900000 86.100000 l 213.600000 86.100000 l 213.300000 85.800000 l 212.700000 85.800000 l 212.400000 86.100000 l 211.800000 86.100000 l 211.500000 85.800000 l 209.700000 85.800000 l 209.400000 85.500000 l 211.500000 85.500000 l 211.800000 85.800000 l 214.200000 85.800000 l 214.500000 85.500000 l 214.200000 85.200000 l 214.200000 84.900000 l 213.900000 84.600000 l 213.300000 84.600000 l 213.000000 84.300000 l 212.100000 84.300000 l 211.800000 84.000000 l 209.400000 84.000000 l 209.100000 84.300000 l 203.700000 84.300000 l 203.400000 84.600000 l 202.800000 84.600000 l 202.500000 84.300000 l 201.000000 84.300000 l 200.400000 84.900000 l 199.800000 84.900000 l 199.500000 84.600000 l 192.600000 84.600000 l 192.300000 84.900000 l 187.800000 84.900000 l 187.500000 84.600000 l 184.800000 84.600000 l 184.500000 84.900000 l 184.200000 84.600000 l 183.600000 85.200000 l 183.000000 85.200000 l 182.700000 84.900000 l 182.400000 85.200000 l 182.100000 84.900000 l 181.200000 84.900000 l 180.900000 85.200000 l 179.400000 85.200000 l 179.100000 84.900000 l 178.200000 84.900000 l 177.900000 85.200000 l 175.800000 85.200000 l 175.500000 85.500000 l 175.200000 85.200000 l 168.000000 85.200000 l 167.700000 85.500000 l 167.700000 85.800000 l 167.400000 85.800000 l 167.100000 86.100000 l 167.100000 86.700000 l 166.800000 87.000000 l 166.800000 88.200000 l 167.100000 88.500000 l 167.400000 88.200000 l 168.600000 88.200000 l 167.700000 88.200000 l 167.400000 88.500000 l 167.100000 88.500000 l 166.800000 88.800000 l 166.800000 89.100000 l 167.100000 89.400000 l 167.400000 89.400000 l 167.700000 89.100000 l 171.900000 89.100000 l 171.900000 89.400000 l 167.700000 89.400000 l 167.400000 89.700000 l h 213.000000 90.300000 m 213.300000 90.600000 l 213.600000 90.600000 l 214.200000 90.000000 l 214.200000 89.700000 l 214.500000 89.400000 l 214.200000 89.100000 l 213.600000 89.100000 l 213.300000 88.800000 l 212.700000 88.800000 l 212.400000 89.100000 l 212.100000 89.100000 l 211.500000 89.700000 l 211.800000 90.000000 l 212.400000 90.000000 l 212.700000 90.300000 l h 206.400000 90.300000 m 206.700000 90.600000 l 207.000000 90.600000 l 207.300000 90.300000 l 207.900000 90.300000 l 208.200000 90.000000 l 209.100000 90.000000 l 209.400000 89.700000 l 209.400000 89.400000 l 209.100000 89.100000 l 208.200000 89.100000 l 207.900000 89.400000 l 207.300000 89.400000 l 207.000000 89.700000 l 206.700000 89.700000 l 206.400000 90.000000 l h 175.800000 90.300000 m 176.100000 90.600000 l 176.700000 90.600000 l 177.600000 89.700000 l 181.200000 89.700000 l 181.500000 90.000000 l 184.200000 90.000000 l 184.500000 89.700000 l 185.400000 89.700000 l 185.700000 89.400000 l 186.300000 89.400000 l 186.600000 89.700000 l 188.100000 89.700000 l 188.400000 90.000000 l 189.300000 90.000000 l 189.600000 89.700000 l 190.200000 89.700000 l 190.500000 89.400000 l 191.400000 89.400000 l 191.700000 89.100000 l 191.400000 88.800000 l 190.500000 88.800000 l 190.200000 88.500000 l 189.900000 88.800000 l 188.100000 88.800000 l 187.800000 88.500000 l 185.400000 88.500000 l 184.800000 87.900000 l 184.500000 87.900000 l 184.200000 87.600000 l 183.600000 88.200000 l 184.200000 88.800000 l 184.800000 88.800000 l 184.500000 88.800000 l 184.200000 89.100000 l 181.500000 89.100000 l 181.500000 88.800000 l 181.200000 88.500000 l 180.900000 88.500000 l 180.600000 88.800000 l 180.600000 89.100000 l 176.100000 89.100000 l 175.800000 89.400000 l 175.500000 89.400000 l 175.200000 89.700000 l 175.800000 90.300000 l 176.400000 89.700000 l h 173.100000 90.300000 m 173.400000 90.600000 l 173.700000 90.300000 l 173.400000 90.000000 l h 166.800000 90.300000 m 167.100000 90.600000 l 167.400000 90.600000 l 167.700000 90.300000 l 167.700000 90.000000 l 167.400000 89.700000 l 167.100000 89.700000 l 166.800000 90.000000 l h 184.500000 90.600000 m 184.800000 90.900000 l 186.600000 90.900000 l 186.900000 90.600000 l 186.600000 90.300000 l 186.300000 90.300000 l 186.000000 90.000000 l 185.400000 90.000000 l 185.100000 90.300000 l 184.800000 90.300000 l h 174.300000 90.600000 m 174.600000 90.900000 l 175.800000 90.900000 l 176.100000 90.600000 l 175.800000 90.300000 l 174.600000 90.300000 l h 184.200000 90.900000 m 184.500000 91.200000 l 184.800000 90.900000 l 184.500000 90.600000 l h 180.600000 91.200000 m 180.900000 91.500000 l 183.900000 91.500000 l 184.200000 91.200000 l 183.900000 90.900000 l 183.300000 90.900000 l 183.000000 90.600000 l 180.900000 90.600000 l 180.600000 90.300000 l 180.300000 90.600000 l 180.000000 90.300000 l 179.700000 90.300000 l 179.400000 90.000000 l 179.100000 90.000000 l 178.500000 90.600000 l 178.800000 90.900000 l 179.100000 90.900000 l 179.400000 91.200000 l h 166.500000 91.200000 m 166.800000 91.500000 l 167.400000 91.500000 l 167.700000 91.200000 l 167.700000 90.900000 l 167.400000 90.600000 l 166.800000 90.600000 l 166.500000 90.900000 l h 166.800000 92.100000 m 167.100000 92.400000 l 167.400000 92.400000 l 167.700000 92.100000 l 167.700000 91.800000 l 167.400000 91.500000 l h 199.200000 92.700000 m 199.500000 93.000000 l 199.800000 93.000000 l 200.100000 92.700000 l 199.800000 92.400000 l 199.500000 92.400000 l h 189.900000 92.700000 m 190.200000 93.000000 l 190.500000 92.700000 l 190.200000 92.400000 l h 212.100000 93.600000 m 212.400000 93.300000 l 214.200000 93.300000 l 214.500000 93.000000 l 214.500000 92.400000 l 214.200000 92.100000 l 213.900000 92.100000 l 213.600000 91.800000 l 212.100000 91.800000 l 211.800000 92.100000 l 209.700000 92.100000 l 209.400000 92.400000 l 209.700000 92.700000 l 211.500000 92.700000 l 211.500000 93.000000 l h 204.000000 93.600000 m 204.300000 93.300000 l 204.900000 93.300000 l 205.200000 93.000000 l 205.500000 93.000000 l 205.800000 92.700000 l 205.500000 92.400000 l 203.700000 92.400000 l 203.400000 92.700000 l 203.400000 93.000000 l h 183.900000 93.300000 m 184.200000 93.600000 l 184.500000 93.300000 l 184.200000 93.000000 l h 204.000000 93.600000 m 204.300000 93.900000 l 204.600000 93.600000 l 204.300000 93.300000 l h 195.600000 93.600000 m 195.900000 93.900000 l 196.800000 93.900000 l 197.100000 93.600000 l 197.400000 93.600000 l 197.700000 93.300000 l 199.200000 93.300000 l 199.500000 93.000000 l 198.900000 92.400000 l 196.200000 92.400000 l 195.900000 92.700000 l 193.200000 92.700000 l 192.900000 92.400000 l 192.600000 92.700000 l 192.300000 92.700000 l 192.000000 93.000000 l 191.100000 93.000000 l 190.800000 93.300000 l 191.100000 93.600000 l 192.600000 93.600000 l 192.900000 93.300000 l 196.500000 93.300000 l 195.900000 93.300000 l h 177.000000 93.900000 m 177.300000 94.200000 l 177.900000 94.200000 l 178.200000 93.900000 l 177.900000 93.600000 l 177.300000 93.600000 l h 204.900000 94.200000 m 205.200000 94.500000 l 205.500000 94.200000 l 205.200000 93.900000 l h 182.400000 94.200000 m 182.700000 94.500000 l 183.600000 94.500000 l 183.900000 94.200000 l 183.900000 93.900000 l 184.200000 93.600000 l 183.900000 93.300000 l 183.600000 93.300000 l 183.000000 93.900000 l 182.700000 93.900000 l h 205.200000 94.500000 m 205.500000 94.800000 l 205.800000 94.500000 l 205.500000 94.200000 l h 168.900000 94.500000 m 169.200000 94.800000 l 169.500000 94.800000 l 169.800000 94.500000 l 170.100000 94.800000 l 170.700000 94.200000 l 170.400000 93.900000 l 168.600000 93.900000 l 168.300000 94.200000 l 168.600000 94.500000 l h 205.500000 94.800000 m 205.800000 95.100000 l 207.600000 95.100000 l 207.900000 94.800000 l 207.600000 94.500000 l 205.800000 94.500000 l h 201.900000 94.800000 m 202.200000 95.100000 l 203.100000 95.100000 l 203.400000 94.800000 l 203.100000 94.500000 l 202.200000 94.500000 l h 167.400000 95.100000 m 168.300000 94.200000 l 168.000000 93.900000 l 167.700000 93.900000 l 167.400000 94.200000 l 167.700000 93.900000 l 167.700000 93.600000 l 172.200000 93.600000 l 172.500000 93.900000 l 172.800000 93.600000 l 175.500000 93.600000 l 175.800000 93.900000 l 176.100000 93.600000 l 176.100000 93.300000 l 176.700000 93.300000 l 177.000000 93.600000 l 177.900000 93.600000 l 178.200000 93.900000 l 178.500000 93.900000 l 178.800000 93.600000 l 179.100000 93.600000 l 179.400000 93.300000 l 181.200000 93.300000 l 181.500000 93.600000 l 181.800000 93.600000 l 180.900000 93.600000 l 180.600000 93.900000 l 180.000000 93.900000 l 179.700000 94.200000 l 180.000000 94.500000 l 181.800000 94.500000 l 183.000000 93.300000 l 183.300000 93.300000 l 183.600000 93.000000 l 184.200000 93.000000 l 184.500000 92.700000 l 185.400000 92.700000 l 185.100000 93.000000 l 184.800000 93.000000 l 184.200000 93.600000 l 184.200000 93.900000 l 184.500000 94.200000 l 184.800000 94.200000 l 186.300000 92.700000 l 186.600000 92.700000 l 186.300000 93.000000 l 186.300000 93.300000 l 186.600000 93.600000 l 187.500000 92.700000 l 187.800000 92.700000 l 188.700000 91.800000 l 189.000000 91.800000 l 189.300000 92.100000 l 188.700000 92.700000 l 189.000000 93.000000 l 189.300000 93.000000 l 190.200000 92.100000 l 190.200000 92.400000 l 190.500000 92.700000 l 190.800000 92.700000 l 191.400000 92.100000 l 191.700000 92.100000 l 192.000000 92.400000 l 192.300000 92.100000 l 192.600000 92.100000 l 192.900000 91.800000 l 194.100000 91.800000 l 194.400000 92.100000 l 199.800000 92.100000 l 200.100000 92.400000 l 201.000000 92.400000 l 201.300000 92.100000 l 202.200000 92.100000 l 202.500000 92.400000 l 202.800000 92.100000 l 204.300000 92.100000 l 204.600000 91.800000 l 205.500000 91.800000 l 205.800000 91.500000 l 206.100000 91.500000 l 206.100000 91.800000 l 205.800000 91.800000 l 205.500000 92.100000 l 205.800000 92.400000 l 206.100000 92.400000 l 205.800000 92.700000 l 206.100000 93.000000 l 206.700000 93.000000 l 207.000000 92.700000 l 207.300000 93.000000 l 209.100000 93.000000 l 209.400000 92.700000 l 209.400000 92.400000 l 209.100000 92.100000 l 207.900000 92.100000 l 207.600000 91.800000 l 207.000000 91.800000 l 208.200000 91.800000 l 208.800000 91.200000 l 208.800000 90.900000 l 208.500000 90.600000 l 207.600000 90.600000 l 207.300000 90.900000 l 205.500000 90.900000 l 204.900000 90.300000 l 204.600000 90.300000 l 204.300000 90.000000 l 201.300000 90.000000 l 201.000000 90.300000 l 200.100000 90.300000 l 199.800000 90.600000 l 199.500000 90.300000 l 196.500000 90.300000 l 196.800000 90.000000 l 196.500000 89.700000 l 196.200000 89.700000 l 195.900000 89.400000 l 195.600000 89.700000 l 194.700000 89.700000 l 194.400000 90.000000 l 194.100000 90.000000 l 193.800000 90.300000 l 193.200000 90.300000 l 192.900000 90.000000 l 192.600000 90.300000 l 192.600000 90.600000 l 189.600000 90.600000 l 189.300000 90.900000 l 189.600000 91.200000 l 190.200000 91.200000 l 190.500000 91.500000 l 190.800000 91.200000 l 191.100000 91.200000 l 191.100000 91.500000 l 190.500000 91.500000 l 190.200000 91.200000 l 189.900000 91.500000 l 189.600000 91.500000 l 189.000000 90.900000 l 188.700000 91.200000 l 188.400000 91.200000 l 188.100000 91.500000 l 188.100000 91.200000 l 187.800000 90.900000 l 187.500000 91.200000 l 186.900000 91.200000 l 186.600000 91.500000 l 186.600000 91.800000 l 186.300000 91.800000 l 186.000000 92.100000 l 185.100000 92.100000 l 184.800000 91.800000 l 184.500000 91.800000 l 184.200000 92.100000 l 183.600000 92.100000 l 183.300000 92.400000 l 183.000000 92.100000 l 182.700000 92.400000 l 180.900000 92.400000 l 180.600000 92.100000 l 180.300000 92.100000 l 180.000000 92.400000 l 176.700000 92.400000 l 177.000000 92.100000 l 178.800000 92.100000 l 179.100000 91.800000 l 178.500000 91.200000 l 176.100000 91.200000 l 175.800000 91.500000 l 175.800000 91.800000 l 175.200000 91.200000 l 174.300000 91.200000 l 174.600000 90.900000 l 174.300000 90.600000 l 173.700000 90.600000 l 173.400000 90.900000 l 173.100000 90.900000 l 172.800000 91.200000 l 172.200000 91.200000 l 172.200000 90.900000 l 172.500000 90.600000 l 172.200000 90.300000 l 167.700000 90.300000 l 167.400000 90.600000 l 167.700000 90.900000 l 170.100000 90.900000 l 170.400000 91.200000 l 171.600000 91.200000 l 167.700000 91.200000 l 167.400000 91.500000 l 167.700000 91.800000 l 168.000000 91.800000 l 168.300000 92.100000 l 168.600000 92.100000 l 167.700000 92.100000 l 166.800000 93.000000 l 166.800000 93.300000 l 167.100000 93.600000 l 166.800000 93.900000 l 166.800000 94.500000 l h 208.500000 95.100000 m 208.800000 95.400000 l 209.400000 95.400000 l 209.700000 95.100000 l 211.200000 95.100000 l 211.500000 94.800000 l 211.800000 95.100000 l 212.100000 95.100000 l 212.400000 95.400000 l 214.200000 95.400000 l 214.500000 95.100000 l 214.500000 94.200000 l 214.200000 93.900000 l 213.900000 93.900000 l 213.600000 94.200000 l 213.900000 94.500000 l 212.400000 94.500000 l 212.400000 94.200000 l 212.100000 93.900000 l 212.100000 93.600000 l 211.500000 93.000000 l 211.200000 93.300000 l 210.300000 93.300000 l 210.000000 93.600000 l 210.300000 93.900000 l 210.000000 94.200000 l 209.700000 94.200000 l 209.400000 94.500000 l 208.800000 94.500000 l 208.500000 94.800000 l h 201.000000 95.700000 m 201.300000 95.400000 l 201.900000 95.400000 l 202.200000 95.100000 l 201.900000 94.800000 l 201.600000 94.800000 l 201.300000 94.500000 l 201.000000 94.800000 l 200.700000 94.500000 l 200.700000 94.200000 l 200.400000 93.900000 l 200.100000 93.900000 l 199.800000 94.200000 l 200.100000 94.500000 l 200.100000 94.800000 l h 196.500000 95.400000 m 196.800000 95.700000 l 197.400000 95.700000 l 197.700000 95.400000 l 198.000000 95.400000 l 198.300000 95.700000 l 198.600000 95.700000 l 199.200000 95.100000 l 198.900000 94.800000 l 198.600000 94.800000 l 198.000000 94.200000 l 197.400000 94.800000 l 197.100000 94.500000 l 196.500000 94.500000 l 196.200000 94.800000 l 196.500000 95.100000 l h 197.400000 95.700000 m 197.700000 96.000000 l 198.000000 96.000000 l 198.300000 95.700000 l 198.000000 95.400000 l 197.700000 95.400000 l h 201.300000 96.900000 m 201.600000 97.200000 l 203.100000 97.200000 l 203.400000 96.900000 l 203.700000 96.900000 l 204.000000 96.600000 l 204.300000 96.600000 l 204.600000 96.300000 l 207.300000 96.300000 l 207.600000 96.000000 l 207.900000 96.000000 l 208.500000 95.400000 l 208.200000 95.100000 l 207.900000 95.100000 l 207.600000 95.400000 l 206.400000 95.400000 l 206.100000 95.700000 l 204.900000 95.700000 l 204.600000 95.400000 l 204.300000 95.700000 l 202.800000 95.700000 l 202.500000 96.000000 l 202.800000 96.300000 l 203.100000 96.300000 l 202.800000 96.600000 l 201.600000 96.600000 l h 167.400000 96.900000 m 167.700000 97.200000 l 168.600000 97.200000 l 168.900000 96.900000 l 168.600000 96.600000 l 167.700000 96.600000 l h 214.200000 97.500000 m 214.500000 97.200000 l 214.500000 96.600000 l 214.200000 96.300000 l 213.300000 96.300000 l 213.000000 96.600000 l 213.300000 96.900000 l 213.600000 96.900000 l h 175.500000 97.200000 m 175.800000 97.500000 l 176.100000 97.200000 l 178.200000 97.200000 l 178.500000 96.900000 l 178.800000 96.900000 l 179.100000 96.600000 l 179.400000 96.600000 l 179.700000 96.900000 l 180.000000 96.600000 l 181.800000 96.600000 l 182.100000 96.300000 l 182.700000 96.300000 l 183.000000 96.600000 l 183.300000 96.300000 l 183.900000 96.300000 l 184.200000 96.600000 l 184.500000 96.300000 l 184.800000 96.300000 l 185.100000 96.000000 l 187.500000 96.000000 l 187.800000 96.300000 l 188.100000 96.300000 l 188.400000 96.600000 l 188.700000 96.600000 l 189.000000 96.300000 l 190.800000 96.300000 l 191.400000 96.900000 l 192.000000 96.300000 l 192.600000 96.300000 l 192.900000 96.000000 l 194.100000 96.000000 l 194.400000 96.300000 l 194.700000 96.300000 l 195.000000 96.000000 l 196.500000 96.000000 l 196.800000 95.700000 l 196.200000 95.100000 l 195.900000 95.100000 l 195.600000 94.800000 l 195.000000 94.800000 l 194.700000 95.100000 l 194.400000 95.100000 l 194.100000 94.800000 l 193.800000 94.800000 l 193.500000 95.100000 l 192.900000 95.100000 l 192.600000 95.400000 l 190.800000 95.400000 l 190.500000 95.700000 l 189.900000 95.700000 l 189.600000 95.400000 l 189.900000 95.400000 l 190.200000 95.100000 l 190.500000 95.100000 l 190.800000 94.800000 l 192.000000 94.800000 l 192.300000 94.500000 l 192.600000 94.500000 l 192.900000 94.200000 l 193.500000 94.200000 l 193.800000 93.900000 l 193.200000 93.300000 l 192.900000 93.300000 l 192.600000 93.600000 l 191.100000 93.600000 l 190.800000 93.900000 l 190.500000 93.600000 l 189.300000 93.600000 l 189.000000 93.900000 l 187.800000 93.900000 l 187.500000 93.600000 l 186.900000 93.600000 l 186.600000 93.900000 l 186.300000 93.600000 l 185.700000 93.600000 l 185.100000 94.200000 l 184.500000 94.200000 l 184.200000 94.500000 l 183.900000 94.500000 l 183.600000 94.800000 l 175.800000 94.800000 l 175.800000 94.500000 l 175.500000 94.200000 l 174.600000 94.200000 l 174.300000 93.900000 l 174.000000 94.200000 l 173.400000 94.200000 l 173.100000 93.900000 l 171.000000 93.900000 l 170.700000 94.200000 l 171.300000 94.800000 l 171.900000 94.800000 l 172.200000 94.500000 l 172.500000 94.500000 l 172.800000 94.800000 l 175.200000 94.800000 l 175.200000 95.100000 l 174.900000 95.100000 l 174.600000 94.800000 l 167.700000 94.800000 l 167.400000 95.100000 l 167.100000 95.100000 l 166.800000 95.400000 l 166.800000 95.700000 l 167.100000 96.000000 l 167.400000 96.000000 l 167.700000 95.700000 l 171.000000 95.700000 l 171.300000 95.400000 l 171.600000 95.700000 l 175.200000 95.700000 l 171.000000 95.700000 l 170.700000 96.000000 l 170.400000 96.000000 l 170.100000 95.700000 l 168.900000 95.700000 l 168.600000 96.000000 l 168.300000 95.700000 l 167.700000 95.700000 l 166.800000 96.600000 l 167.100000 96.900000 l 167.400000 96.900000 l 167.700000 96.600000 l 172.200000 96.600000 l 172.500000 96.900000 l 172.800000 96.600000 l 174.600000 96.600000 l 174.900000 96.900000 l 175.500000 96.900000 l h 175.800000 96.000000 m 176.100000 95.700000 l 177.000000 95.700000 l 177.300000 95.400000 l 177.600000 95.700000 l 182.400000 95.700000 l 176.100000 95.700000 l h 166.800000 97.200000 m 167.100000 97.500000 l 167.400000 97.500000 l 167.700000 97.200000 l 167.400000 96.900000 l 167.100000 96.900000 l h 189.900000 97.800000 m 190.200000 98.100000 l 190.500000 98.100000 l 190.800000 97.800000 l 192.600000 97.800000 l 192.900000 97.500000 l 192.600000 97.200000 l 192.300000 97.200000 l 191.700000 96.600000 l 191.400000 96.900000 l 191.400000 97.200000 l 188.700000 97.200000 l 188.400000 97.500000 l 188.700000 97.800000 l h 189.600000 98.400000 m 189.900000 98.700000 l 190.200000 98.400000 l 189.900000 98.100000 l h 190.500000 99.300000 m 190.800000 99.300000 l 191.100000 99.000000 l 190.500000 98.400000 l 190.200000 98.400000 l 189.900000 98.700000 l h 167.100000 99.000000 m 167.400000 99.300000 l 167.700000 99.000000 l 167.700000 98.700000 l 168.000000 98.700000 l 168.600000 98.100000 l 168.900000 98.100000 l 169.200000 97.800000 l 168.900000 97.500000 l 167.700000 97.500000 l 167.100000 98.100000 l h 190.800000 99.300000 m 191.100000 99.600000 l 191.400000 99.300000 l 191.100000 99.000000 l h 213.600000 102.000000 m 213.900000 101.700000 l 213.900000 100.500000 l 213.600000 100.200000 l 213.300000 100.200000 l 213.000000 100.500000 l 213.300000 100.800000 l 213.300000 101.100000 l 213.000000 101.400000 l h 167.400000 102.300000 m 167.700000 102.600000 l 168.000000 102.300000 l 168.000000 101.400000 l 167.700000 101.100000 l 167.100000 101.700000 l 167.400000 102.000000 l h 167.100000 102.900000 m 167.400000 103.200000 l 167.700000 102.900000 l 167.700000 102.600000 l 167.400000 102.300000 l 167.100000 102.600000 l h 212.700000 107.400000 m 213.000000 107.700000 l 213.300000 107.700000 l 213.600000 107.400000 l 213.600000 107.100000 l 213.900000 106.800000 l 213.900000 104.100000 l 213.600000 103.800000 l 213.900000 103.500000 l 213.900000 103.200000 l 213.600000 102.900000 l 213.600000 102.000000 l 213.300000 101.700000 l 213.000000 102.000000 l 213.000000 103.500000 l 212.700000 103.800000 l 213.000000 104.100000 l 212.700000 104.400000 l 213.000000 104.700000 l 212.700000 105.000000 l 213.000000 105.300000 l 212.700000 105.600000 l 212.700000 105.900000 l 213.000000 106.200000 l 212.700000 106.500000 l h 190.500000 107.700000 m 190.800000 108.000000 l 191.100000 107.700000 l 190.800000 107.400000 l h 213.000000 108.000000 m 213.300000 108.300000 l 213.600000 108.300000 l 213.900000 108.000000 l 213.900000 107.700000 l 213.600000 107.400000 l h 212.400000 108.000000 m 212.700000 108.300000 l 213.000000 108.000000 l 212.700000 107.700000 l h 190.800000 108.000000 m 191.100000 108.300000 l 191.400000 108.000000 l 191.100000 107.700000 l h 193.500000 116.700000 m 193.800000 117.000000 l 194.100000 116.700000 l 193.800000 116.400000 l h 196.500000 117.000000 m 196.800000 117.300000 l 197.100000 117.000000 l 198.000000 117.000000 l 198.600000 116.400000 l 198.600000 115.200000 l 198.900000 114.900000 l 198.900000 114.600000 l 198.600000 114.300000 l 198.900000 114.000000 l 198.600000 113.700000 l 198.900000 113.400000 l 198.900000 111.900000 l 198.600000 111.600000 l 198.600000 110.400000 l 198.900000 110.100000 l 198.900000 109.800000 l 200.100000 108.600000 l 200.700000 109.200000 l 201.000000 108.900000 l 201.600000 108.900000 l 201.900000 108.600000 l 202.200000 108.900000 l 202.500000 108.600000 l 203.100000 108.600000 l 203.400000 108.300000 l 204.600000 108.300000 l 204.900000 108.600000 l 205.500000 108.600000 l 205.800000 108.300000 l 206.100000 108.300000 l 206.400000 108.000000 l 206.400000 107.700000 l 206.100000 107.400000 l 206.100000 106.800000 l 205.800000 106.500000 l 205.800000 105.900000 l 206.100000 105.600000 l 206.100000 105.000000 l 205.800000 104.700000 l 206.100000 104.400000 l 206.100000 104.100000 l 205.800000 103.800000 l 205.800000 103.500000 l 206.100000 103.200000 l 206.100000 102.900000 l 205.800000 102.600000 l 205.800000 102.000000 l 206.400000 101.400000 l 206.400000 101.100000 l 206.700000 100.800000 l 206.700000 99.900000 l 207.000000 99.900000 l 206.700000 100.200000 l 206.700000 100.800000 l 207.000000 101.100000 l 207.000000 101.700000 l 206.700000 102.000000 l 207.000000 102.300000 l 206.400000 102.900000 l 206.700000 103.200000 l 206.400000 103.500000 l 206.700000 103.800000 l 206.700000 107.400000 l 207.000000 107.700000 l 207.000000 108.000000 l 207.600000 108.600000 l 207.900000 108.300000 l 208.200000 108.600000 l 208.500000 108.600000 l 209.100000 108.000000 l 209.100000 107.700000 l 208.800000 107.400000 l 208.800000 106.200000 l 208.500000 105.900000 l 208.500000 103.200000 l 208.800000 102.900000 l 209.100000 102.900000 l 209.100000 103.500000 l 209.400000 103.800000 l 209.100000 104.100000 l 209.100000 105.300000 l 209.400000 105.600000 l 209.400000 106.800000 l 209.100000 107.100000 l 209.100000 107.400000 l 209.400000 107.700000 l 209.700000 107.700000 l 209.700000 108.000000 l 210.000000 108.300000 l 210.300000 108.000000 l 210.900000 108.000000 l 211.200000 108.300000 l 211.500000 108.000000 l 212.400000 108.000000 l 212.700000 107.700000 l 212.400000 107.400000 l 212.400000 105.600000 l 212.700000 105.300000 l 212.400000 105.000000 l 212.400000 104.100000 l 212.100000 103.800000 l 212.100000 103.200000 l 212.400000 102.900000 l 212.400000 102.000000 l 212.700000 101.700000 l 212.700000 101.400000 l 212.400000 101.100000 l 212.400000 100.800000 l 212.700000 100.500000 l 212.400000 100.200000 l 212.400000 99.300000 l 212.100000 99.000000 l 212.100000 98.700000 l 212.400000 98.400000 l 212.700000 98.400000 l 212.700000 99.000000 l 213.300000 99.600000 l 213.300000 99.900000 l 213.600000 100.200000 l 213.900000 99.900000 l 213.900000 99.000000 l 213.600000 98.700000 l 213.900000 98.400000 l 213.900000 98.100000 l 213.600000 97.800000 l 213.300000 97.800000 l 213.000000 97.500000 l 212.400000 97.500000 l 211.800000 96.900000 l 210.600000 96.900000 l 210.300000 96.600000 l 210.000000 96.900000 l 209.700000 96.600000 l 209.400000 96.900000 l 209.400000 97.200000 l 207.600000 97.200000 l 207.300000 96.900000 l 207.000000 97.200000 l 204.600000 97.200000 l 204.300000 97.500000 l 204.000000 97.500000 l 203.700000 97.200000 l 203.400000 97.500000 l 201.300000 97.500000 l 201.000000 97.800000 l 200.100000 97.800000 l 199.800000 97.500000 l 195.000000 97.500000 l 194.700000 97.800000 l 194.400000 97.800000 l 194.400000 97.500000 l 194.700000 97.200000 l 194.400000 96.900000 l 192.900000 96.900000 l 192.600000 97.200000 l 192.900000 97.500000 l 193.500000 97.500000 l 193.200000 97.800000 l 193.200000 98.100000 l 192.900000 98.400000 l 192.900000 99.300000 l 193.200000 99.600000 l 193.200000 100.200000 l 194.100000 101.100000 l 194.400000 101.100000 l 195.300000 102.000000 l 195.300000 107.700000 l 195.600000 108.000000 l 195.600000 109.800000 l 195.300000 110.100000 l 195.600000 110.400000 l 195.600000 113.100000 l 195.300000 113.400000 l 195.300000 116.100000 l 195.600000 116.400000 l 195.600000 116.700000 l 195.900000 117.000000 l h 208.500000 102.000000 m 208.500000 100.800000 l 208.800000 100.500000 l 208.800000 100.200000 l 208.500000 99.900000 l 208.800000 99.600000 l 208.800000 99.300000 l 208.500000 99.000000 l 209.100000 99.000000 l 209.400000 99.300000 l 209.400000 101.700000 l 209.100000 102.000000 l h 210.900000 107.400000 m 210.600000 107.100000 l 210.900000 106.800000 l 210.900000 106.200000 l 210.600000 105.900000 l 210.600000 105.300000 l 210.900000 105.000000 l 210.900000 104.700000 l 210.600000 104.400000 l 210.600000 102.000000 l 210.900000 101.700000 l 210.900000 101.400000 l 211.200000 101.100000 l 210.900000 100.800000 l 210.900000 99.300000 l 210.600000 99.000000 l 210.600000 98.100000 l 210.900000 98.100000 l 211.200000 98.400000 l 211.200000 99.000000 l 211.500000 99.300000 l 211.500000 100.800000 l 211.800000 101.100000 l 211.500000 101.400000 l 211.500000 101.700000 l 211.200000 102.000000 l 211.200000 102.900000 l 210.900000 103.200000 l 211.200000 103.500000 l 211.200000 107.400000 l h 184.500000 117.300000 m 184.800000 117.000000 l 185.700000 117.000000 l 186.000000 117.300000 l 186.300000 117.000000 l 186.900000 117.000000 l 187.200000 117.300000 l 187.500000 117.000000 l 187.500000 116.400000 l 187.200000 116.100000 l 186.900000 116.100000 l 187.200000 115.800000 l 187.200000 115.200000 l 186.900000 114.900000 l 187.500000 114.300000 l 187.200000 114.000000 l 187.500000 113.700000 l 187.500000 112.500000 l 187.800000 112.200000 l 187.500000 111.900000 l 187.500000 108.600000 l 187.200000 108.300000 l 187.500000 108.000000 l 187.200000 107.700000 l 187.200000 106.500000 l 186.900000 106.200000 l 187.200000 105.900000 l 187.200000 104.700000 l 186.900000 104.400000 l 187.200000 104.100000 l 187.200000 103.800000 l 186.900000 103.500000 l 187.200000 103.200000 l 187.200000 102.600000 l 187.500000 102.300000 l 187.500000 102.000000 l 187.200000 101.700000 l 187.200000 101.400000 l 187.500000 101.400000 l 187.800000 101.100000 l 187.800000 100.800000 l 188.100000 101.100000 l 188.700000 100.500000 l 188.700000 99.600000 l 188.400000 99.300000 l 188.400000 98.400000 l 188.100000 98.100000 l 187.800000 98.100000 l 187.500000 97.800000 l 184.500000 97.800000 l 184.200000 98.100000 l 183.300000 98.100000 l 183.000000 98.400000 l 179.700000 98.400000 l 179.400000 98.100000 l 179.100000 98.400000 l 178.800000 98.100000 l 178.200000 98.100000 l 177.900000 98.400000 l 177.600000 98.100000 l 177.300000 98.100000 l 177.000000 98.400000 l 176.100000 98.400000 l 175.800000 98.700000 l 175.500000 98.400000 l 169.500000 98.400000 l 169.200000 98.700000 l 168.900000 98.700000 l 168.600000 99.000000 l 168.900000 99.300000 l 168.900000 99.600000 l 169.200000 99.900000 l 169.500000 99.900000 l 169.500000 101.100000 l 169.800000 101.400000 l 169.500000 101.700000 l 169.500000 102.300000 l 169.200000 102.300000 l 169.200000 100.800000 l 168.900000 100.500000 l 168.900000 99.900000 l 168.600000 99.600000 l 168.900000 99.300000 l 168.600000 99.000000 l 168.300000 99.000000 l 167.400000 99.900000 l 167.400000 100.200000 l 167.700000 100.500000 l 168.300000 100.500000 l 168.000000 100.800000 l 168.300000 101.100000 l 168.300000 101.400000 l 168.000000 101.700000 l 168.000000 102.300000 l 168.300000 102.600000 l 168.000000 102.900000 l 168.000000 103.200000 l 167.700000 103.500000 l 168.000000 103.800000 l 168.000000 105.600000 l 167.700000 105.900000 l 168.000000 106.200000 l 168.000000 107.400000 l 167.700000 107.700000 l 167.700000 108.000000 l 168.300000 108.600000 l 168.600000 108.600000 l 168.900000 108.900000 l 169.200000 108.600000 l 169.200000 108.300000 l 168.900000 108.000000 l 168.900000 106.500000 l 169.200000 106.500000 l 169.200000 106.800000 l 168.900000 107.100000 l 168.900000 107.700000 l 169.200000 108.000000 l 169.200000 108.900000 l 169.500000 109.200000 l 170.400000 109.200000 l 170.700000 108.900000 l 170.700000 108.000000 l 170.400000 107.700000 l 170.400000 106.500000 l 170.700000 106.200000 l 170.700000 106.500000 l 171.000000 106.800000 l 170.700000 107.100000 l 170.700000 107.400000 l 171.000000 107.700000 l 171.000000 108.900000 l 171.300000 109.200000 l 172.200000 109.200000 l 172.500000 109.500000 l 173.100000 109.500000 l 173.400000 109.200000 l 173.700000 109.500000 l 174.000000 109.500000 l 174.300000 109.200000 l 174.600000 109.500000 l 174.900000 109.200000 l 175.200000 109.500000 l 175.500000 109.200000 l 175.800000 109.200000 l 176.100000 109.500000 l 177.300000 109.500000 l 177.600000 109.200000 l 180.600000 109.200000 l 180.900000 109.500000 l 182.400000 109.500000 l 182.700000 109.800000 l 183.000000 109.800000 l 183.600000 110.400000 l 183.600000 111.000000 l 183.900000 111.300000 l 183.900000 112.200000 l 184.200000 112.500000 l 184.200000 113.400000 l 183.900000 113.700000 l 184.200000 114.000000 l 183.900000 114.300000 l 183.900000 114.600000 l 183.600000 114.900000 l 183.900000 115.200000 l 183.900000 115.500000 l 183.600000 115.800000 l 183.600000 116.100000 l 183.900000 116.400000 l 183.900000 116.700000 l h 170.400000 105.600000 m 170.400000 104.700000 l 170.100000 104.400000 l 170.400000 104.100000 l 170.400000 102.900000 l 170.700000 102.900000 l 171.000000 102.600000 l 170.700000 102.300000 l 170.700000 101.400000 l 170.400000 101.100000 l 170.400000 99.600000 l 170.700000 99.900000 l 171.000000 99.900000 l 171.300000 100.200000 l 171.000000 100.500000 l 171.000000 101.700000 l 171.300000 102.000000 l 171.000000 102.300000 l 171.300000 102.600000 l 171.000000 102.600000 l 170.700000 102.900000 l 171.000000 103.200000 l 170.700000 103.500000 l 170.700000 105.600000 l h 172.500000 108.600000 m 172.500000 108.300000 l 172.200000 108.000000 l 172.200000 106.500000 l 171.900000 106.200000 l 172.200000 105.900000 l 172.200000 105.000000 l 171.900000 104.700000 l 171.900000 104.400000 l 172.200000 104.100000 l 172.200000 103.800000 l 171.900000 103.500000 l 171.900000 102.900000 l 172.200000 102.900000 l 172.500000 102.600000 l 172.500000 101.100000 l 172.800000 100.800000 l 172.500000 100.500000 l 172.800000 100.800000 l 172.800000 102.600000 l 172.500000 102.900000 l h 188.700000 117.300000 m 189.000000 117.600000 l 190.500000 117.600000 l 190.800000 117.300000 l 191.100000 117.600000 l 191.400000 117.300000 l 191.700000 117.300000 l 192.000000 117.000000 l 192.000000 116.400000 l 191.400000 115.800000 l 191.700000 115.500000 l 191.400000 115.200000 l 191.400000 114.900000 l 192.000000 114.300000 l 191.700000 114.000000 l 191.700000 113.700000 l 192.000000 113.400000 l 192.000000 112.200000 l 191.700000 111.900000 l 191.700000 111.600000 l 192.000000 111.300000 l 191.700000 111.000000 l 192.000000 110.700000 l 192.000000 109.800000 l 192.300000 109.500000 l 192.300000 109.200000 l 191.400000 108.300000 l 190.800000 108.900000 l 190.800000 112.500000 l 191.100000 112.800000 l 191.100000 113.100000 l 190.800000 113.400000 l 190.800000 113.100000 l 190.500000 112.800000 l 190.500000 110.100000 l 190.800000 109.800000 l 190.500000 109.500000 l 190.800000 109.200000 l 190.500000 108.900000 l 190.500000 108.000000 l 189.900000 107.400000 l 190.200000 107.100000 l 190.200000 106.800000 l 190.500000 106.800000 l 190.200000 107.100000 l 190.500000 107.400000 l 190.800000 107.400000 l 191.100000 107.700000 l 191.400000 107.400000 l 191.400000 106.800000 l 191.700000 106.500000 l 191.400000 106.200000 l 191.400000 105.900000 l 191.700000 105.600000 l 191.400000 105.300000 l 191.400000 105.000000 l 191.700000 104.700000 l 191.700000 104.400000 l 191.400000 104.100000 l 191.400000 102.900000 l 192.000000 102.300000 l 191.700000 102.600000 l 191.700000 103.200000 l 192.000000 103.500000 l 192.000000 104.100000 l 191.700000 104.400000 l 192.000000 104.700000 l 192.000000 107.700000 l 192.300000 108.000000 l 192.000000 108.300000 l 192.300000 108.600000 l 192.300000 111.900000 l 192.000000 112.200000 l 192.300000 112.500000 l 192.300000 114.300000 l 192.000000 114.600000 l 192.000000 117.000000 l 192.600000 117.600000 l 192.900000 117.600000 l 193.200000 117.300000 l 193.500000 117.600000 l 193.800000 117.300000 l 193.800000 117.000000 l 193.500000 116.700000 l 193.200000 116.700000 l 193.200000 116.400000 l 192.900000 116.100000 l 192.900000 114.900000 l 193.200000 114.600000 l 193.200000 113.100000 l 193.500000 112.800000 l 193.200000 112.500000 l 193.200000 109.800000 l 192.900000 109.500000 l 193.200000 109.200000 l 193.200000 108.000000 l 192.900000 107.700000 l 192.900000 105.600000 l 192.600000 105.300000 l 192.600000 105.000000 l 192.900000 104.700000 l 192.600000 104.400000 l 192.900000 104.100000 l 192.600000 103.800000 l 192.900000 103.500000 l 192.600000 103.200000 l 192.600000 102.900000 l 192.900000 102.600000 l 192.900000 102.300000 l 193.200000 102.000000 l 193.200000 101.400000 l 192.900000 101.100000 l 192.600000 101.400000 l 189.000000 101.400000 l 188.700000 101.700000 l 188.400000 101.700000 l 187.500000 102.600000 l 187.500000 103.200000 l 187.800000 103.500000 l 187.800000 103.800000 l 187.500000 104.100000 l 187.500000 104.400000 l 187.800000 104.700000 l 187.800000 106.800000 l 187.500000 107.100000 l 187.500000 107.400000 l 188.100000 108.000000 l 188.100000 110.400000 l 187.800000 110.700000 l 188.100000 111.000000 l 188.100000 112.500000 l 187.800000 112.800000 l 188.100000 113.100000 l 188.100000 114.300000 l 187.800000 114.600000 l 187.800000 116.100000 l 188.100000 116.400000 l 188.100000 117.000000 l 188.400000 117.300000 l h 190.200000 116.700000 m 190.500000 116.400000 l 190.200000 116.100000 l 190.500000 115.800000 l 190.200000 115.500000 l 190.200000 114.600000 l 190.500000 114.600000 l 190.800000 114.300000 l 190.800000 114.000000 l 190.800000 114.600000 l 190.500000 114.900000 l 190.500000 115.200000 l 190.800000 115.500000 l 190.800000 115.800000 l 190.500000 116.100000 l 190.800000 116.400000 l 190.800000 116.700000 l h 192.300000 157.200000 m 192.300000 156.900000 l 192.000000 156.600000 l 192.000000 155.700000 l 191.700000 155.400000 l 191.700000 155.100000 l 191.100000 154.500000 l 191.100000 153.300000 l 190.800000 153.000000 l 190.500000 153.000000 l 190.200000 152.700000 l 190.200000 151.500000 l 189.600000 150.900000 l 189.300000 151.200000 l 189.300000 150.600000 l 189.000000 150.300000 l 189.000000 150.000000 l 188.700000 149.700000 l 188.700000 149.400000 l 188.400000 149.100000 l 188.400000 148.500000 l 187.800000 147.900000 l 187.500000 147.900000 l 187.200000 147.600000 l 187.200000 147.000000 l 186.900000 146.700000 l 186.900000 146.100000 l 186.300000 145.500000 l 186.000000 145.800000 l 185.700000 145.500000 l 185.700000 144.900000 l 185.400000 144.600000 l 185.400000 144.300000 l 185.100000 144.000000 l 184.800000 144.000000 l 184.500000 143.700000 l 184.200000 143.700000 l 184.200000 143.400000 l 183.900000 143.100000 l 183.900000 142.800000 l 183.600000 142.500000 l 183.600000 141.600000 l 183.300000 141.300000 l 182.700000 141.300000 l 182.400000 141.000000 l 182.400000 140.400000 l 182.100000 140.100000 l 182.100000 139.500000 l 181.800000 139.200000 l 181.800000 138.900000 l 182.100000 138.600000 l 180.900000 137.400000 l 180.900000 137.100000 l 180.600000 136.800000 l 180.600000 136.200000 l 180.900000 135.900000 l 180.600000 135.600000 l 180.600000 135.300000 l 180.300000 135.000000 l 180.300000 134.100000 l 180.000000 133.800000 l 180.000000 133.500000 l 179.400000 132.900000 l 179.700000 132.600000 l 179.700000 132.300000 l 179.400000 132.000000 l 179.400000 131.700000 l 179.100000 131.400000 l 179.100000 129.300000 l 179.400000 129.000000 l 179.400000 127.800000 l 179.700000 127.500000 l 179.700000 126.900000 l 180.000000 126.600000 l 180.000000 125.400000 l 180.300000 125.100000 l 180.000000 124.800000 l 180.000000 124.500000 l 180.600000 123.900000 l 180.600000 123.600000 l 180.900000 123.600000 l 181.200000 123.300000 l 181.200000 123.000000 l 182.100000 122.100000 l 182.100000 121.800000 l 184.800000 119.100000 l 184.800000 118.800000 l 188.100000 118.800000 l 188.400000 119.100000 l 192.900000 119.100000 l 193.200000 118.800000 l 195.900000 118.800000 l 196.200000 118.500000 l 196.500000 118.800000 l 198.000000 118.800000 l 198.600000 119.400000 l 198.600000 119.700000 l 198.900000 120.000000 l 198.900000 120.300000 l 198.600000 120.600000 l 198.300000 120.600000 l 198.000000 120.900000 l 198.000000 121.200000 l 198.600000 121.800000 l 199.800000 121.800000 l 200.400000 122.400000 l 200.400000 123.300000 l 200.700000 123.600000 l 201.000000 123.600000 l 201.000000 123.900000 l 201.900000 124.800000 l 201.900000 125.100000 l 202.200000 125.400000 l 202.200000 126.000000 l 201.900000 126.300000 l 202.200000 126.600000 l 202.200000 126.900000 l 202.500000 127.200000 l 202.500000 128.100000 l 202.800000 128.400000 l 202.800000 129.000000 l 202.500000 129.300000 l 203.100000 129.900000 l 203.100000 130.500000 l 202.800000 130.800000 l 202.800000 131.400000 l 202.500000 131.700000 l 202.500000 132.600000 l 202.800000 132.900000 l 202.500000 133.200000 l 202.500000 133.800000 l 202.200000 134.100000 l 202.500000 134.400000 l 202.500000 134.700000 l 202.200000 135.000000 l 201.900000 135.000000 l 201.600000 135.300000 l 201.600000 135.600000 l 201.300000 135.900000 l 201.300000 136.800000 l 201.600000 137.100000 l 201.300000 137.400000 l 201.300000 138.300000 l 201.000000 138.600000 l 200.700000 138.600000 l 200.400000 138.900000 l 200.700000 139.200000 l 200.100000 139.800000 l 200.100000 140.400000 l 199.800000 140.400000 l 199.200000 141.000000 l 199.500000 141.300000 l 199.500000 141.900000 l 199.200000 142.200000 l 199.200000 142.500000 l 198.900000 142.800000 l 198.600000 142.500000 l 198.300000 142.800000 l 198.300000 143.100000 l 198.000000 143.400000 l 198.000000 144.000000 l 197.700000 144.300000 l 197.700000 144.900000 l 197.400000 145.200000 l 197.400000 145.500000 l 196.200000 146.700000 l 196.200000 147.300000 l 195.900000 147.600000 l 195.600000 147.600000 l 195.300000 147.900000 l 195.300000 149.400000 l 195.000000 149.700000 l 195.300000 150.000000 l 194.100000 151.200000 l 194.100000 152.400000 l 193.500000 153.000000 l 193.500000 153.300000 l 193.200000 153.600000 l 193.200000 155.100000 l 193.500000 155.400000 l 193.200000 155.700000 l 193.200000 156.300000 l 192.900000 156.600000 l 192.900000 157.200000 l h 197.700000 120.600000 m 198.000000 120.900000 l 198.300000 120.600000 l 198.000000 120.300000 l h 197.400000 121.800000 m 197.700000 122.100000 l 198.000000 121.800000 l 197.700000 121.500000 l h 199.200000 124.800000 m 199.500000 124.500000 l 199.500000 123.900000 l 199.200000 123.600000 l 199.200000 123.300000 l 198.900000 123.000000 l 198.900000 122.700000 l 198.000000 121.800000 l 197.700000 122.100000 l 197.700000 122.400000 l 198.300000 123.000000 l 198.300000 123.600000 l 198.600000 123.900000 l 198.600000 124.200000 l h 182.100000 124.800000 m 182.400000 125.100000 l 182.700000 124.800000 l 182.400000 124.500000 l h 199.800000 126.300000 m 200.100000 126.600000 l 200.400000 126.300000 l 200.400000 125.400000 l 200.100000 125.100000 l 200.100000 124.800000 l 199.800000 124.500000 l 199.500000 124.800000 l 199.500000 125.400000 l 199.800000 125.700000 l h 181.800000 126.600000 m 182.100000 126.900000 l 182.400000 126.900000 l 182.700000 126.600000 l 182.700000 125.700000 l 183.000000 125.400000 l 183.000000 125.100000 l 182.700000 124.800000 l 182.400000 125.100000 l 182.100000 125.100000 l 181.800000 125.400000 l h 181.200000 127.800000 m 181.500000 128.100000 l 181.800000 127.800000 l 181.800000 127.500000 l 182.100000 127.200000 l 182.100000 126.900000 l 181.800000 126.600000 l 181.200000 127.200000 l h 180.900000 128.100000 m 181.200000 128.400000 l 181.500000 128.100000 l 181.200000 127.800000 l h 200.100000 128.700000 m 200.400000 129.000000 l 200.700000 128.700000 l 200.700000 128.100000 l 200.400000 127.800000 l 200.100000 128.100000 l h 200.400000 129.600000 m 200.700000 129.900000 l 201.000000 129.600000 l 201.000000 129.000000 l 200.700000 128.700000 l 200.400000 129.000000 l h 199.800000 132.300000 m 200.100000 132.000000 l 200.100000 131.700000 l 200.400000 131.400000 l 200.100000 131.100000 l 200.100000 130.800000 l 199.800000 130.500000 l 199.800000 128.400000 l 199.500000 128.100000 l 199.500000 127.500000 l 199.200000 127.200000 l 199.500000 126.900000 l 199.500000 125.700000 l 198.900000 125.100000 l 198.900000 124.800000 l 198.600000 124.500000 l 198.600000 124.200000 l 198.300000 123.900000 l 198.300000 123.600000 l 197.700000 123.000000 l 197.700000 122.700000 l 197.100000 122.100000 l 196.800000 122.400000 l 196.800000 122.700000 l 197.100000 123.000000 l 197.100000 123.300000 l 197.400000 123.600000 l 197.400000 124.200000 l 198.300000 125.100000 l 198.300000 125.700000 l 198.600000 126.000000 l 198.600000 126.300000 l 198.300000 126.600000 l 198.300000 126.900000 l 198.600000 127.200000 l 198.600000 128.100000 l 198.900000 128.400000 l 198.900000 129.900000 l 199.200000 130.200000 l 199.200000 130.800000 l 199.500000 131.100000 l 199.200000 131.400000 l 199.200000 131.700000 l h 183.600000 138.000000 m 183.900000 138.300000 l 184.200000 138.000000 l 184.200000 135.900000 l 183.900000 135.600000 l 183.600000 135.600000 l 183.300000 135.900000 l 183.300000 137.100000 l 183.600000 137.400000 l h 197.700000 139.200000 m 198.000000 139.500000 l 198.300000 139.200000 l 198.000000 138.900000 l h 197.400000 140.100000 m 197.700000 140.400000 l 198.000000 140.400000 l 198.300000 140.100000 l 198.300000 139.800000 l 197.700000 139.200000 l 197.400000 139.500000 l h 183.600000 140.700000 m 183.900000 141.000000 l 184.200000 140.700000 l 184.200000 139.200000 l 183.900000 138.900000 l 183.600000 138.900000 l 183.300000 139.200000 l 183.600000 139.500000 l h 196.200000 141.300000 m 196.500000 141.600000 l 197.100000 141.000000 l 196.800000 140.700000 l 196.800000 140.100000 l 197.100000 139.800000 l 197.100000 138.600000 l 196.800000 138.300000 l 196.200000 138.900000 l h 185.400000 141.600000 m 185.700000 141.300000 l 185.700000 141.000000 l 185.400000 140.700000 l 185.400000 140.400000 l 185.700000 140.100000 l 185.700000 139.800000 l 185.100000 139.200000 l 184.800000 139.500000 l 184.800000 141.000000 l h 185.400000 141.600000 m 185.700000 141.900000 l 186.000000 141.600000 l 185.700000 141.300000 l h 196.200000 143.100000 m 196.500000 143.400000 l 196.800000 143.100000 l 196.800000 142.800000 l 197.100000 142.500000 l 197.100000 141.600000 l 196.800000 141.300000 l 196.500000 141.600000 l 196.500000 142.200000 l 196.200000 142.500000 l h 186.900000 144.900000 m 187.200000 145.200000 l 187.500000 144.900000 l 187.200000 144.600000 l h 188.400000 145.500000 m 188.700000 145.800000 l 189.000000 145.500000 l 188.700000 145.200000 l h 193.800000 145.800000 m 194.100000 146.100000 l 194.400000 145.800000 l 194.100000 145.500000 l h 188.700000 147.000000 m 189.000000 147.300000 l 189.300000 147.000000 l 189.300000 145.800000 l 189.000000 145.500000 l 188.400000 146.100000 l 188.700000 146.400000 l h 189.900000 148.200000 m 190.200000 148.500000 l 190.500000 148.200000 l 190.500000 147.600000 l 190.800000 147.300000 l 190.500000 147.000000 l 190.500000 146.700000 l 190.200000 146.400000 l 189.900000 146.400000 l 189.600000 146.700000 l 189.900000 147.000000 l h 192.600000 148.500000 m 192.900000 148.800000 l 193.200000 148.500000 l 193.200000 147.900000 l 192.900000 147.600000 l 192.900000 147.300000 l 192.600000 147.000000 l 192.300000 147.300000 l 192.300000 147.900000 l 192.600000 148.200000 l h 191.100000 148.500000 m 191.400000 148.800000 l 191.700000 148.500000 l 191.700000 147.900000 l 192.000000 147.600000 l 192.000000 147.300000 l 191.700000 147.000000 l 192.000000 146.700000 l 192.000000 146.100000 l 191.700000 145.800000 l 192.000000 145.500000 l 191.700000 145.200000 l 192.000000 144.900000 l 192.000000 144.300000 l 191.700000 144.000000 l 192.000000 143.700000 l 191.700000 143.400000 l 192.000000 143.400000 l 192.300000 143.700000 l 192.300000 144.000000 l 192.000000 144.300000 l 192.000000 144.900000 l 192.300000 145.200000 l 192.300000 146.700000 l 192.900000 147.300000 l 193.200000 147.000000 l 193.200000 146.700000 l 192.900000 146.400000 l 192.900000 145.800000 l 193.200000 145.500000 l 193.200000 143.400000 l 192.900000 143.100000 l 193.200000 142.800000 l 193.200000 141.600000 l 193.500000 141.600000 l 193.500000 144.000000 l 193.200000 144.300000 l 193.500000 144.600000 l 193.500000 144.900000 l 193.200000 145.200000 l 193.200000 145.500000 l 193.500000 145.800000 l 193.800000 145.800000 l 194.100000 145.500000 l 194.100000 144.900000 l 194.400000 144.600000 l 194.400000 143.400000 l 194.100000 143.100000 l 194.400000 142.800000 l 194.400000 141.900000 l 194.700000 141.600000 l 194.100000 141.000000 l 194.400000 140.700000 l 194.100000 140.400000 l 194.400000 140.100000 l 194.400000 138.900000 l 194.100000 138.600000 l 194.400000 138.600000 l 194.700000 138.300000 l 194.700000 138.000000 l 194.400000 137.700000 l 194.400000 137.400000 l 194.700000 137.100000 l 194.700000 136.800000 l 194.400000 136.500000 l 194.700000 136.200000 l 194.700000 135.300000 l 195.000000 135.000000 l 194.700000 134.700000 l 194.700000 133.200000 l 194.400000 132.900000 l 194.100000 132.900000 l 194.100000 132.600000 l 194.400000 132.600000 l 195.000000 133.200000 l 195.000000 133.800000 l 195.300000 134.100000 l 195.000000 134.400000 l 195.300000 134.700000 l 195.300000 135.000000 l 195.000000 135.300000 l 195.000000 135.600000 l 195.300000 135.900000 l 195.300000 137.100000 l 195.000000 137.400000 l 195.300000 137.700000 l 195.600000 137.700000 l 194.700000 138.600000 l 194.700000 140.100000 l 194.400000 140.400000 l 194.700000 140.700000 l 194.700000 141.300000 l 195.000000 141.600000 l 195.000000 142.500000 l 194.700000 142.800000 l 195.000000 143.100000 l 194.700000 143.400000 l 195.300000 144.000000 l 195.600000 143.700000 l 195.600000 143.400000 l 195.900000 143.100000 l 195.600000 142.800000 l 195.900000 142.500000 l 195.900000 139.500000 l 196.200000 139.200000 l 195.900000 138.900000 l 195.900000 138.600000 l 196.200000 138.600000 l 196.500000 138.300000 l 196.500000 137.700000 l 196.200000 137.400000 l 196.500000 137.100000 l 196.200000 136.800000 l 196.200000 135.900000 l 196.500000 135.600000 l 196.200000 135.300000 l 196.500000 135.000000 l 196.200000 134.700000 l 196.200000 134.400000 l 196.500000 134.100000 l 196.200000 133.800000 l 195.900000 133.800000 l 196.200000 133.500000 l 196.200000 133.200000 l 195.900000 132.900000 l 196.200000 132.600000 l 196.200000 132.300000 l 195.900000 132.000000 l 195.900000 131.100000 l 195.600000 130.800000 l 195.900000 130.500000 l 195.900000 129.900000 l 195.900000 130.200000 l 196.200000 130.500000 l 196.200000 132.900000 l 196.800000 133.500000 l 196.800000 134.100000 l 196.500000 134.400000 l 196.800000 134.700000 l 196.500000 135.000000 l 196.800000 135.300000 l 196.500000 135.600000 l 196.800000 135.900000 l 196.800000 136.500000 l 196.500000 136.800000 l 196.800000 137.100000 l 196.800000 138.300000 l 197.100000 138.600000 l 197.400000 138.600000 l 197.700000 138.300000 l 197.700000 133.200000 l 197.400000 132.900000 l 197.400000 132.000000 l 197.100000 131.700000 l 197.100000 131.400000 l 197.400000 131.100000 l 197.400000 130.500000 l 197.100000 130.200000 l 197.100000 129.000000 l 196.800000 128.700000 l 196.800000 128.100000 l 196.500000 127.800000 l 196.500000 126.600000 l 196.800000 126.300000 l 196.500000 126.000000 l 196.500000 124.800000 l 196.800000 125.100000 l 196.800000 126.000000 l 197.100000 126.300000 l 196.800000 126.600000 l 196.800000 126.900000 l 197.100000 127.200000 l 197.100000 128.400000 l 197.400000 128.700000 l 197.400000 129.600000 l 197.700000 129.900000 l 197.700000 130.500000 l 198.000000 130.800000 l 197.700000 131.100000 l 198.000000 131.400000 l 198.000000 132.000000 l 197.700000 132.300000 l 198.000000 132.600000 l 197.700000 132.900000 l 198.000000 133.200000 l 198.300000 133.200000 l 198.300000 134.100000 l 198.000000 134.400000 l 198.000000 135.000000 l 198.300000 135.300000 l 198.000000 135.600000 l 198.300000 135.900000 l 198.000000 136.200000 l 198.300000 136.500000 l 198.300000 136.800000 l 198.600000 137.100000 l 198.900000 136.800000 l 198.900000 136.200000 l 199.800000 135.300000 l 199.500000 135.000000 l 199.500000 134.700000 l 199.200000 134.400000 l 199.200000 133.800000 l 198.900000 133.500000 l 198.900000 133.200000 l 199.200000 132.900000 l 198.900000 132.600000 l 198.600000 132.600000 l 198.600000 132.300000 l 198.900000 132.000000 l 198.600000 131.700000 l 198.900000 131.400000 l 198.900000 130.200000 l 198.600000 129.900000 l 198.600000 128.400000 l 198.000000 127.800000 l 198.000000 126.900000 l 198.300000 126.600000 l 198.300000 126.300000 l 198.000000 126.000000 l 198.000000 125.700000 l 197.700000 125.400000 l 197.700000 125.100000 l 197.100000 124.500000 l 197.100000 123.900000 l 196.800000 123.600000 l 196.800000 123.000000 l 196.500000 122.700000 l 196.500000 122.400000 l 195.900000 121.800000 l 195.900000 121.200000 l 195.600000 120.900000 l 195.900000 120.900000 l 196.500000 120.300000 l 196.500000 120.600000 l 196.800000 120.900000 l 197.100000 120.900000 l 198.000000 120.000000 l 197.700000 119.700000 l 193.500000 119.700000 l 192.900000 120.300000 l 188.400000 120.300000 l 187.800000 119.700000 l 187.500000 120.000000 l 185.400000 120.000000 l 184.500000 120.900000 l 184.200000 120.900000 l 183.000000 122.100000 l 183.000000 122.400000 l 182.700000 122.700000 l 182.700000 123.000000 l 182.400000 123.300000 l 182.400000 124.500000 l 182.700000 124.800000 l 183.000000 124.500000 l 183.000000 124.200000 l 183.600000 123.600000 l 183.600000 123.300000 l 183.600000 124.200000 l 183.300000 124.500000 l 183.600000 124.800000 l 183.600000 125.100000 l 183.300000 125.400000 l 183.300000 125.700000 l 183.000000 126.000000 l 183.000000 126.600000 l 182.400000 127.200000 l 182.400000 128.400000 l 182.100000 128.700000 l 182.100000 129.900000 l 181.800000 130.200000 l 181.800000 131.100000 l 182.100000 131.400000 l 181.800000 131.700000 l 181.800000 132.600000 l 182.100000 132.900000 l 182.400000 132.900000 l 182.700000 132.600000 l 183.000000 132.600000 l 183.000000 132.900000 l 183.300000 133.200000 l 183.300000 133.500000 l 183.600000 133.800000 l 184.200000 133.800000 l 183.600000 133.800000 l 183.300000 134.100000 l 183.600000 134.400000 l 183.300000 134.700000 l 183.600000 135.000000 l 183.900000 135.000000 l 184.200000 134.700000 l 184.200000 135.000000 l 184.500000 135.300000 l 184.200000 135.600000 l 184.500000 135.900000 l 184.500000 136.800000 l 184.800000 137.100000 l 184.500000 137.400000 l 184.500000 137.700000 l 184.800000 138.000000 l 184.800000 138.600000 l 184.500000 138.900000 l 184.500000 139.200000 l 184.800000 139.500000 l 185.400000 138.900000 l 185.400000 138.600000 l 185.700000 138.300000 l 185.400000 138.000000 l 185.700000 137.700000 l 185.700000 137.400000 l 186.300000 138.000000 l 186.000000 138.300000 l 186.300000 138.600000 l 186.000000 138.600000 l 185.700000 138.900000 l 186.000000 139.200000 l 186.000000 141.000000 l 186.600000 141.600000 l 186.600000 142.800000 l 186.900000 143.100000 l 186.900000 144.300000 l 187.500000 144.900000 l 187.800000 144.600000 l 187.800000 144.300000 l 187.500000 144.000000 l 187.500000 143.700000 l 187.800000 143.400000 l 187.500000 143.100000 l 187.500000 141.600000 l 187.200000 141.300000 l 187.500000 141.300000 l 187.800000 141.600000 l 187.800000 141.900000 l 188.100000 142.200000 l 188.100000 142.500000 l 187.800000 142.800000 l 188.100000 143.100000 l 188.100000 144.000000 l 188.400000 144.300000 l 188.400000 144.900000 l 189.000000 145.500000 l 189.600000 144.900000 l 189.300000 144.600000 l 189.600000 144.900000 l 189.600000 146.100000 l 189.900000 146.400000 l 190.200000 146.400000 l 190.500000 146.100000 l 190.500000 145.200000 l 190.800000 144.900000 l 190.500000 144.600000 l 190.500000 142.200000 l 190.200000 141.900000 l 190.200000 141.600000 l 189.900000 141.300000 l 189.900000 140.700000 l 190.200000 140.400000 l 190.200000 140.100000 l 189.900000 139.800000 l 189.900000 138.900000 l 190.200000 138.600000 l 190.200000 138.300000 l 190.500000 138.300000 l 190.500000 138.600000 l 190.200000 138.900000 l 190.200000 139.500000 l 190.500000 139.800000 l 190.200000 140.100000 l 190.200000 140.700000 l 190.500000 141.000000 l 190.200000 141.300000 l 190.500000 141.600000 l 190.800000 141.600000 l 190.800000 142.500000 l 190.500000 142.800000 l 190.500000 143.100000 l 190.800000 143.400000 l 190.800000 145.500000 l 191.100000 145.800000 l 190.800000 146.100000 l 190.800000 146.400000 l 191.100000 146.700000 l h 192.000000 142.500000 m 191.700000 142.200000 l 192.000000 141.900000 l 192.000000 141.600000 l 191.700000 141.300000 l 191.400000 141.300000 l 191.400000 140.100000 l 191.100000 139.800000 l 191.100000 139.500000 l 191.400000 139.200000 l 191.400000 138.900000 l 191.700000 138.900000 l 192.000000 138.600000 l 191.700000 138.300000 l 191.700000 136.800000 l 192.000000 137.100000 l 192.000000 137.400000 l 192.300000 137.700000 l 192.300000 138.000000 l 192.000000 138.300000 l 192.000000 138.600000 l 191.700000 138.900000 l 191.700000 139.500000 l 192.000000 139.800000 l 192.000000 140.100000 l 191.700000 140.400000 l 192.000000 140.700000 l 191.700000 141.000000 l 191.700000 141.300000 l 192.000000 141.600000 l 192.300000 141.600000 l 192.000000 141.900000 l 192.300000 142.200000 l 192.300000 142.500000 l h f* Q}\special{pdf:econtent}}
\put(0,0){\special{pdf:bcontent}\special{pdf:code q 0 g 21.300000 13.800000 m 21.000000 13.500000 l 21.000000 12.300000 l 21.300000 12.000000 l 21.300000 11.700000 l 21.900000 11.100000 l 22.200000 11.400000 l 23.400000 11.400000 l 24.000000 10.800000 l 24.300000 10.800000 l 24.600000 11.100000 l 28.500000 11.100000 l 28.800000 11.400000 l 31.800000 11.400000 l 32.100000 11.100000 l 33.300000 11.100000 l 33.600000 10.800000 l 33.900000 11.100000 l 40.200000 11.100000 l 40.500000 10.800000 l 40.800000 11.100000 l 45.900000 11.100000 l 46.200000 11.400000 l 46.500000 11.100000 l 46.800000 11.400000 l 48.300000 11.400000 l 48.600000 11.100000 l 49.200000 11.100000 l 49.500000 10.800000 l 49.800000 11.100000 l 50.100000 10.800000 l 50.400000 11.100000 l 50.700000 11.100000 l 51.000000 10.800000 l 51.600000 10.800000 l 51.900000 11.100000 l 52.200000 11.100000 l 52.500000 10.800000 l 53.100000 11.400000 l 53.700000 11.400000 l 54.000000 11.100000 l 54.300000 11.400000 l 55.500000 11.400000 l 55.800000 11.100000 l 56.100000 11.400000 l 56.700000 11.400000 l 57.000000 11.100000 l 57.300000 11.100000 l 57.600000 10.800000 l 58.200000 10.800000 l 58.500000 11.100000 l 58.800000 10.800000 l 59.700000 10.800000 l 60.000000 11.100000 l 60.300000 11.100000 l 60.600000 10.800000 l 60.900000 11.100000 l 61.200000 10.800000 l 62.400000 10.800000 l 62.700000 11.100000 l 63.900000 11.100000 l 64.200000 10.800000 l 64.800000 10.800000 l 65.100000 11.100000 l 65.400000 10.800000 l 66.600000 10.800000 l 66.900000 11.100000 l 67.200000 11.100000 l 67.500000 10.800000 l 67.800000 11.100000 l 68.100000 11.100000 l 68.400000 10.800000 l 68.700000 10.800000 l 69.000000 11.100000 l 69.600000 11.100000 l 69.900000 10.800000 l 70.200000 10.800000 l 70.500000 11.100000 l 71.100000 11.100000 l 71.400000 10.800000 l 71.700000 11.100000 l 72.000000 10.800000 l 72.300000 11.100000 l 73.200000 11.100000 l 73.500000 10.800000 l 73.800000 10.800000 l 74.100000 10.500000 l 74.400000 10.800000 l 74.700000 10.500000 l 75.300000 10.500000 l 75.600000 10.800000 l 75.900000 10.500000 l 76.200000 10.500000 l 76.500000 10.800000 l 76.800000 10.500000 l 77.100000 10.800000 l 77.400000 10.800000 l 77.700000 11.100000 l 78.000000 10.800000 l 79.800000 10.800000 l 80.100000 11.100000 l 80.400000 10.800000 l 80.700000 11.100000 l 81.900000 11.100000 l 82.500000 10.500000 l 82.800000 10.800000 l 84.300000 10.800000 l 84.600000 10.500000 l 84.900000 10.500000 l 85.200000 10.800000 l 86.400000 10.800000 l 86.700000 10.500000 l 87.000000 10.800000 l 87.300000 10.500000 l 87.600000 10.800000 l 87.900000 10.500000 l 88.800000 10.500000 l 89.100000 10.800000 l 90.000000 10.800000 l 90.300000 11.100000 l 90.600000 10.800000 l 90.600000 10.500000 l 90.900000 10.800000 l 91.200000 10.500000 l 91.500000 10.800000 l 92.400000 10.800000 l 93.000000 10.200000 l 93.300000 10.500000 l 94.800000 10.500000 l 95.100000 10.800000 l 95.400000 10.500000 l 98.700000 10.500000 l 99.300000 9.900000 l 99.600000 10.200000 l 101.100000 10.200000 l 101.400000 10.500000 l 104.700000 10.500000 l 105.000000 10.800000 l 105.000000 12.000000 l 104.700000 12.300000 l 99.300000 12.300000 l 99.000000 12.000000 l 98.700000 12.300000 l 98.700000 12.600000 l 98.400000 12.300000 l 98.100000 12.600000 l 97.800000 12.300000 l 96.900000 12.300000 l 96.600000 12.600000 l 96.000000 12.600000 l 95.700000 12.300000 l 95.100000 12.300000 l 94.800000 12.600000 l 90.600000 12.600000 l 90.300000 12.900000 l 89.100000 12.900000 l 88.800000 12.600000 l 82.500000 12.600000 l 82.200000 12.900000 l 77.100000 12.900000 l 76.800000 12.600000 l 74.100000 12.600000 l 73.800000 12.900000 l 65.700000 12.900000 l 65.400000 13.200000 l 65.100000 13.200000 l 64.800000 12.900000 l 57.300000 12.900000 l 57.000000 13.200000 l 52.800000 13.200000 l 52.500000 12.900000 l 48.900000 12.900000 l 48.600000 13.200000 l 40.800000 13.200000 l 40.500000 12.900000 l 40.200000 13.200000 l 31.800000 13.200000 l 31.500000 13.500000 l 29.400000 13.500000 l 29.100000 13.200000 l 28.800000 13.500000 l 28.500000 13.200000 l 23.700000 13.200000 l 23.400000 13.500000 l 22.200000 13.500000 l 21.900000 13.800000 l h f* Q}\special{pdf:econtent}}
\put(62.4,186.9){\makebox(0,0){\fontsize{15bp}{17bp}\selectfont A}}
\put(61.2,64.2){\makebox(0,0){\fontsize{15bp}{17bp}\selectfont B}}
\put(13.2,104.7){\makebox(0,0){\rotatebox[origin=c]{90}{\fontsize{15bp}{17bp}\selectfont Ocular}}}
\put(57,129.3){\makebox(0,0){\rotatebox[origin=c]{45}{\fontsize{15bp}{17bp}\selectfont Spiegel}}}
\put(125.1,11.7){\makebox(0,0){\fontsize{15bp}{17bp}\selectfont Fäden}}
\end{picture}
\endgroup

\end{center}

Durch Versuche überzeugte ich mich bald, dass jenes undeutliche Lichtflammenbild blos durch Reflexion der Lichtflamme von der Linse B entsteht. Dies Bild trägt aber gar nichts dazu bei, die Fäden oder deren von Quecksilber reflectirtes Bild sichtbar zu machen, sondern erschwert dies nur in hohem Grade. Ich nehme also die Linse B ganz heraus. Das Ocular vergrössert dann allerdings etwas weniger, aber immer noch genug, und gibt herrliche scharfe Bilder. Das Ocular z. B. welches ich am Meridiankreise anwende, vergrössert mit A und B 134 mal, mit A allein 96 mal.

Ich habe so, mag ich Lampe oder Tageslicht (letzteres etwa nur vermittelst eines in freier Hand gehaltenen Spiegelstücks) gebrauchen, ein kleines, aber immer noch genügendes, gleichförmig beleuchtetes Feld. Der wesentliche Punkt ist, dass zwischen dem unter $45^\circ$ geneigten Spiegel und dem Fadensystem kein Glas sein darf.
